%vim:foldmethod=marker
\documentclass{article}
\usepackage{natbib}
\usepackage{paper}
\usepackage{xcolor}
\usepackage{tikz}


% ----------------------------- Title & Author ------------------------------ %
\title{Causal Inference with Ordinal Outcomes: Flexible Estimation Framework Using Normalizing Flows}
\author{Chanhyuk Park}
\date{}
% --------------------------------------------------------------------------- %

\begin{document}
    \maketitle

\begin{abstract}
  Ordinal outcomes, such as Likert-type survey responses, are ubiquitous in political science, yet they pose distinctive challenges for causal inference. These outcomes are routinely analyzed by assigning numerical scores and applying linear models, despite longstanding concerns that ordinal responses do not identify treatment effects on an underlying cardinal scale without additional assumptions. Building on the classical identification result for ordinal outcomes and existing recommendations to report probability-based quantities, I develop a flexible likelihood-based framework for estimating distributional causal effects, including category-specific and cumulative treatment effects, which are invariant to admissible transformations of the latent scale. Existing approaches typically estimate these quantities using parametric ordered logit or ordered probit models, which rely on restrictive assumptions about the latent error distribution and index specification.
  This paper proposes a flexible likelihood-based estimation framework for ordinal causal inference using conditional normalizing flows. The proposed estimator flexibly learns the conditional distribution underlying observed ordinal responses without imposing a parametric latent error distribution, while preserving the likelihood-based structure familiar from classical ordered response models. The same framework accommodates both structured latent-index models and fully nonparametric conditional distribution estimators, allowing researchers to balance interpretability and modeling flexibility within a unified estimation procedure.
  Monte Carlo simulations show that the proposed estimator maintains low bias across a broad range of latent error distributions and nonlinear response mechanisms under which conventional ordered response models perform poorly. Replications of two published survey experiments demonstrate that flexible distributional estimation can materially alter substantive conclusions about experimental treatment effects.
\end{abstract}

\newpage
\section{Introduction}
\label{sec:introduction}

Causal inference has become a central organizing principle in political science. Experimental and quasi-experimental designs, including survey experiments, difference-in-differences, and regression discontinuity designs, are now standard empirical strategies for studying the effects of policies, institutions, and information. Survey experiments are particularly prominent because random assignment provides a transparent basis for causal identification. Between 2021 and 2024, roughly 20\% of articles published in the \textit{American Political Science Review}, \textit{American Journal of Political Science}, and \textit{Journal of Politics} employed a survey experiment as part of their empirical strategy.

Many of these studies investigate outcomes that are inherently latent, including support for redistribution \citep{Alt2017a, Magni2021a}, evaluations of political leaders \citep{Canes-Wrone2002a, Kriner2009a}, and foreign policy preferences \citep{Tomz2020a}. Although these concepts are commonly viewed as lying on an underlying continuous attitude scale, they are almost always observed through ordinal response categories such as Likert-type scales ranging from \textit{strongly disagree} to \textit{strongly agree}. Consequently, causal inference in political science routinely relies on outcomes that preserve the ordering of responses but not the distances between them.

Applied researchers typically analyze ordinal outcomes in one of three ways. The most common approach assigns numerical scores to response categories and estimates treatment effects using difference-in-means or ordinary least squares, implicitly treating ordinal responses as cardinal. Another common strategy collapses multiple categories into a binary outcome before estimation, sacrificing information about intermediate responses. A third approach fits ordered response models, such as ordered logit or ordered probit, which explicitly account for the ordinal nature of the data through a latent-variable representation. Each approach reflects a different set of assumptions about the relationship between the observed ordinal responses and the underlying latent attitude.

This paper begins from a classical observation in the ordinal response literature: latent treatment effects are identified only up to location and scale transformations of the underlying latent variable. Consequently, treatment effects defined through arbitrary numerical codings of ordinal responses generally require additional measurement assumptions linking the observed categories to a cardinal latent scale \citep{Schroder2017a, Bond2018a, Bloem2022a}. This observation motivates a focus on distributional causal quantities, such as category-specific and cumulative treatment effects, which are defined directly on the observed ordinal responses and do not rely on assumptions about the cardinal spacing between adjacent categories.

This perspective complements recent recommendations in political methodology to summarize ordinal treatment effects using predicted probabilities rather than latent coefficients \citep{Hanmer2013a}. Rather than revisiting how ordered response models should be interpreted, this paper asks how these probability-based causal quantities should be estimated when the assumptions underlying conventional ordered response models are difficult to justify. The methodological focus therefore shifts from interpreting latent coefficients to flexibly estimating the conditional ordinal response distribution. Existing semiparametric estimators relax some of these assumptions but remain difficult to apply with rich covariate structures or nonlinear response mechanisms that are common in modern survey research.

Estimating these distributional causal effects requires modeling the conditional distribution of ordinal responses. Ordered logit and ordered probit remain the workhorses of applied political science because they efficiently exploit the ordering of response categories, admit interpretable latent-variable representations, and often perform well when their assumptions are approximately satisfied. Their validity, however, depends on assumptions about the latent error distribution and index specification that may be difficult to justify in many survey experiments. Political attitudes are frequently polarized or clustered, producing multimodal latent distributions, while persuasive treatments often affect moderate respondents differently from strong supporters, generating nonlinear treatment-response relationships. In such settings, misspecification can distort estimated category probabilities and cumulative treatment effects, even when the experimental design itself identifies the causal effect. The practical challenge is therefore not causal identification, but flexible estimation of the conditional ordinal response distribution.

This paper develops a flexible likelihood-based framework for estimating causal effects with ordinal outcomes. The proposed estimator employs conditional normalizing flows to model the conditional response distribution without imposing a fixed parametric latent error distribution. Unlike existing semiparametric estimators, the framework remains computationally feasible with high-dimensional covariates while preserving exact likelihood-based inference. A single estimation procedure spans classical latent-index models, generalized latent-index specifications, and fully nonparametric conditional response distributions, allowing researchers to balance structural assumptions and modeling flexibility while targeting the same probability-based causal quantities.

I evaluate the proposed framework through Monte Carlo simulations designed to reflect empirically realistic survey settings, including polarized latent attitudes, nonlinear treatment-response relationships, and non-Gaussian latent error distributions. Across these settings, the proposed estimator maintains low bias while conventional ordered response models exhibit substantial distortions in estimated probability effects. I then revisit the survey experiments of \citet{Tomz2020a} and \citet{Mattingly2025a}, showing that flexible estimation of the ordinal response distribution can materially alter substantive conclusions regarding the magnitude and distribution of treatment effects.

The remainder of the paper proceeds as follows. Section~\ref{sec:causalordinal} formalizes causal inference with ordinal outcomes and motivates distributional causal estimands. Section~\ref{sec:estimation} reviews existing estimation approaches and introduces the proposed likelihood-based framework. Section~\ref{sec:simulation} evaluates finite-sample performance through Monte Carlo simulations. Section~\ref{sec:application} presents the empirical applications. Section~\ref{sec:conclusion} concludes.

\clearpage



\section{Causal Inference with Ordinal Outcomes}
\label{sec:causalordinal}

This section develops the conceptual framework underlying the remainder of the paper. I begin by describing ordinal survey responses as measurements of an underlying latent attitude using the classical threshold-crossing model. I then revisit a well-known identification result showing that latent treatment effects are identified only up to location and scale, and discuss its implications for causal inference with ordinal outcomes. Finally, I introduce distributional causal estimands that are directly identified from observed ordinal responses and motivate the estimation framework developed in the next section.

\subsection{Latent Outcomes, Ordinal Responses, and Causal Effects}
\label{subsec:setup}

% fig:latent {{{
\begin{figure}[ht]
\centering
\begin{tikzpicture}[>=stealth,scale=1]

% Axis
\draw[->,thick] (-5,0) -- (5.5,0) node[right] {$Y_i^\ast$};

% Thresholds
\foreach \x/\lab in {-3/$\alpha_0$,-1/$\alpha_1$,1/$\alpha_2$,3/$\alpha_3$}
{
    \draw[thick] (\x,-0.18) -- (\x,0.18);
    \node[below] at (\x,-0.18) {\lab};
}

% Categories
\node at (-4,0.55) {\small SD};
\node at (-2,0.55) {\small D};
\node at (0,0.55) {\small N};
\node at (2,0.55) {\small A};
\node at (4,0.55) {\small SA};

% Latent distribution
\draw[thick,blue]
plot[smooth]
coordinates{
(-4.5,0)
(-3.8,.25)
(-3,.55)
(-2,.95)
(-1,1.3)
(0,1.55)
(1,1.35)
(2,.95)
(3,.45)
(4,.12)
};

% Treatment shift
\draw[thick,red,dashed]
plot[smooth]
coordinates{
(-3.7,0)
(-3,.18)
(-2,.42)
(-1,.78)
(0,1.15)
(1,1.55)
(2,1.45)
(3,1.0)
(4,.45)
(5,.1)
};

\node[blue] at (-1.2,1.7) {\small Control};
\node[red] at (2.0,1.75) {\small Treatment};

\end{tikzpicture}

\caption{Ordinal responses arise by partitioning an underlying latent attitude into ordered response categories. The observed data identify how treatment changes probabilities of falling into response categories, but not the absolute scale of the latent attitude.}
\label{fig:latent}
\end{figure}
% }}}

Many quantities of interest in political science, such as support for redistribution \citep{Alt2017a, Magni2021a}, immigration preferences \citep{Mayda2005a}, trust in institutions, and foreign policy attitudes \citep{Scheve2001r, Tomz2020a}, are naturally viewed as continuous latent constructs. Survey respondents, however, rarely report these attitudes on a continuous scale. Instead, researchers observe ordinal responses, such as Likert items ranging from \textit{strongly disagree} to \textit{strongly agree}, ordered approval categories, or frequency scales. Consequently, empirical analyses are conducted using outcomes that preserve the ordering of respondents’ attitudes but not the distances between adjacent response categories.

Following the classical ordered response framework, let $Y_i^\ast$ denote respondent $i$’s latent attitude and let $Y_i\in{0,1,\ldots,J}$ denote the observed ordinal response. The observed response is generated through a threshold-crossing mechanism,

\[
Y_i=j \quad\Longleftrightarrow\quad \alpha_{j-1}<Y_i^\ast\le\alpha_j,
\qquad j=0,\ldots,J,
\]

where

\[
-\infty=\alpha_{-1}<\alpha_0<\cdots<\alpha_J=\infty
\]

are unknown thresholds partitioning the latent scale. Figure~\ref{fig:latent} illustrates this measurement process. Throughout the paper I assume respondents share a common reporting function. Section~XX discusses extensions allowing reporting thresholds to vary across respondents following generalized ordered response models \citep{Williams2006a}. Alternative survey designs, such as anchored vignettes \citep{King2004a, King2007a}, provide complementary approaches for addressing heterogeneous reporting behavior.

The threshold-crossing representation should be interpreted as a measurement model rather than a structural model. It formalizes the idea that ordinal responses preserve the ordering of the underlying attitudes while discarding information about the distances between adjacent categories. Consequently, the observed responses provide only partial information about the latent construct: they reveal which interval of the latent scale a respondent occupies, but not the respondent’s exact location within that interval.

To describe causal effects, let $Y_i^\ast(d)$ denote the latent potential outcome under treatment level $d$, and let

\[
Y_i(d)=g\left(Y_i^\ast(d)\right)
\]

denote the corresponding observed ordinal potential outcome. Throughout the paper I adopt the standard assumptions of SUTVA, conditional ignorability, and positivity \citep{Rubin1974a}. If the latent outcome $Y_i^\ast$ were directly observed, these assumptions would suffice to identify conventional causal estimands such as

\[
ATE^\ast = \mathbb{E} = \left[ Y_i^\ast(1)-Y_i^\ast(0) \right].
\]

In practice, however, researchers observe only the ordinal response $Y_i$. The central question is therefore not whether causal effects are identified, but which causal quantities remain identified once the underlying latent attitude is observed only through ordered response categories. The next subsection revisits the classical identification result for ordinal response models and uses it to distinguish between causal quantities that require additional measurement assumptions and those that do not.

\subsection{Identification of Causal Effects from Ordinal Outcomes}
\label{subsec:identification}

The threshold-crossing framework provides a useful representation of ordinal outcomes, but it also reveals a fundamental limitation of what can be learned from ordinal data. A classical result in the ordered response literature is that latent-variable models are identified only up to location and scale. Although this normalization is standard in econometrics and psychometrics, its implications for causal inference have received comparatively little attention in applied political science. In particular, survey experiments frequently interpret treatment effects estimated from numerically coded ordinal responses as effects on an underlying latent attitude without making explicit the measurement assumptions required for such an interpretation.

To formalize this point, suppose the latent outcome satisfies

\begin{equation}
  Y_i^\ast = h_{\theta}(X_i) + D_i^\top\tau + \varepsilon_i,
\label{eq:latent}
\end{equation}

where $h(X_i,\beta)$ denotes the systematic component of the latent outcome, $\tau$ is the treatment effect on the latent scale, and $\varepsilon_i$ has cumulative distribution function $F$. Combining this latent-index specification with the threshold-crossing measurement model implies

\[
\begin{aligned}
  P(Y_i=j\mid X_i,D_i) &= F\left( \alpha_j-h(X_i,\beta)-D_i^\top\tau \right) \\
   &= F\left( \alpha_{j-1}-h(X_i,\beta)-D_i^\top\tau \right). \\
\end{aligned}
\]

Now consider any constants $a\in\mathbb{R}$ and $c>0$, and define a transformed latent outcome

\[
\tilde Y_i^\ast = a+cY_i^\ast,
\]

together with transformed thresholds

\[
\tilde\alpha_j = a+c\alpha_j.
\]

Because both the latent outcome and the thresholds are transformed simultaneously, the probabilities of the observed ordinal responses remain unchanged.

\begin{theorem}[Location and Scale Nonidentification]
Under the threshold-crossing latent-variable model, any positive affine transformation of the latent outcome and thresholds generates the same distribution of the observed ordinal responses. Consequently, treatment effects on the latent outcome are identified only up to location and scale.
\end{theorem}

Theorem~1 should not be interpreted as a failure of causal identification. Under the standard assumptions of SUTVA, conditional ignorability, and positivity, causal effects remain identified for any well-defined outcome. Rather, the theorem highlights a measurement problem. Ordinal responses preserve the ordering of latent attitudes but do not reveal the cardinal scale on which those attitudes are measured. Consequently, causal quantities defined on the latent scale require additional measurement assumptions beyond those needed for causal identification itself.

This distinction has immediate implications for empirical practice. Different analyses of ordinal outcomes target different causal quantities and therefore rely on different measurement assumptions. For example, assigning numerical scores to ordinal categories and estimating mean differences implicitly assumes that those scores provide a meaningful cardinal representation of the latent attitude. By contrast, causal quantities defined directly on the observed ordinal responses do not require assumptions about the cardinal spacing between adjacent categories. The next subsection formalizes these probability-based causal estimands and shows how they arise naturally from the observed ordinal response distribution.


\subsection{Distributional Causal Estimands}
\label{subsec:estimands}

The identification result in Theorem~1 does not imply that causal inference with ordinal outcomes is fundamentally limited. Rather, it suggests that inference should focus on causal quantities defined directly on the observed ordinal responses. Such quantities require no assumptions about the cardinal spacing of the latent attitude and therefore remain invariant to admissible transformations of the latent scale. This perspective is closely aligned with recommendations to report predicted probabilities rather than latent coefficients when interpreting ordinal response models \citep{Hanmer2013a}. Here I formalize these quantities as causal estimands and show that they are identified under the standard assumptions of causal inference.

The fundamental object is the distribution of potential ordinal outcomes. Let

\[
\pi_j(d) = P\left( Y_i(d)=j \right),
\qquad j=0,\ldots,J,
\]

denote the probability that a randomly selected individual reports category $j$ under treatment level $d$.

Under the standard assumptions of SUTVA, conditional ignorability, and positivity, these probabilities are identified by standardization,

\[
\pi_j(d) = \mathbb{E}_X\left[P\left( Y_i=j \mid D_i=d,X_i \right)\right],
\]

where the expectation is taken over the covariate distribution of the target population. In randomized experiments, this expression simplifies to the observed proportion of respondents falling into category $j$ within each treatment arm.

\begin{proposition}
Under SUTVA, conditional ignorability, and positivity, the distribution of potential ordinal outcomes is identified by standardization. Consequently, any causal quantity that is a functional of this distribution is also identified.
\end{proposition}

This proposition has an important practical implication. Once the conditional ordinal response distribution has been estimated, a wide range of causal quantities become available without fitting additional models. The primary statistical task is therefore estimation of the conditional response probabilities,

\[
P(Y_i=j\mid D_i,X_i),
\]

which serve as the building blocks for all distributional causal estimands considered in this paper.

The simplest causal quantity is the category-specific treatment effect,

\[
\Delta_j = \pi_j(1)-\pi_j(0),
\]

which measures how treatment changes the probability that a respondent selects category $j$. For example, in a survey experiment on support for military intervention, $\Delta_{\text{Strongly Agree}}$ measures the change in the proportion of respondents expressing the strongest level of support.

In many political science applications, however, interest centers on whether respondents exceed a substantively meaningful threshold rather than on a single response category. This motivates cumulative treatment effects,

\[
\Delta_{\ge j} = P\left( Y_i(1)\ge j \right) = P\left( Y_i(0)\ge j \right),
\]

which quantify changes in the probability of responding at or above category $j$. For example, a researcher may wish to estimate whether an informational treatment increases the probability that respondents express at least moderate support for a policy. Unlike dichotomizing the outcome before estimation, cumulative treatment effects can be calculated for every threshold after estimating the ordinal response distribution, allowing the full information contained in the ordinal responses to be retained.

These estimands also clarify the relationship between the proposed framework and existing empirical practice. Assigning numerical scores to ordinal categories estimates

\[
\mathbb{E} = \left[ s(Y_i(1))-s(Y_i(0)) \right],
\]

where $s(\cdot)$ denotes the chosen scoring rule. This quantity is well defined for a fixed coding rule but generally does not identify a treatment effect on the underlying latent attitude. Similarly, dichotomizing an ordinal outcome estimates one cumulative treatment effect corresponding to a single threshold while discarding information about the remaining response categories. Estimating the full ordinal response distribution instead makes both category-specific and cumulative treatment effects available simultaneously, allowing researchers to select the quantity most appropriate for their substantive question after estimation rather than before it.

The discussion above highlights an important feature of distributional causal inference. Rather than targeting a single summary measure, the primary object of interest is the conditional distribution of ordinal responses,
\[
P(Y_i=j\mid D_i,X_i),
\]
from which a wide range of causal quantities can be obtained. Category-specific treatment effects and cumulative treatment effects are the primary estimands considered in this paper because they are directly interpretable and commonly reported in applied political science. More generally, however, any functional of the estimated response distributions may be calculated after estimation, including global measures of distributional change such as the one-dimensional Wasserstein (earth mover's) distance. Consequently, the statistical problem addressed in the remainder of this paper is estimation of the conditional ordinal response distribution rather than any single causal summary.


\section{Flexible Estimation of Ordinal Response Distributions}
\label{sec:estimation}

Section~\ref{sec:causalordinal} established that the causal quantities of interest are functionals of the conditional ordinal response distribution,
\[
P(Y_i=j\mid D_i,X_i).
\]
The remaining statistical problem is therefore estimation of this distribution. This section first reviews existing estimation approaches and their underlying assumptions before introducing a flexible likelihood-based framework that generalizes classical ordered response models while preserving their probabilistic interpretation.

\subsection{Existing Estimation Approaches}
\label{subsec:existing}

The dominant approach to modeling ordinal outcomes in political science is the ordered response model, most commonly ordered logit or ordered probit. These models combine the latent-index representation in Equation~\eqref{eq:latent} with the threshold-crossing measurement model introduced in Section~\ref{subsec:setup}. The distinction between ordered probit and ordered logit lies only in the assumed distribution of the latent error term, with the former assuming a standard normal distribution and the latter assuming a standard logistic distribution.

Ordered response models remain the workhorses of applied political science because they naturally respect the ordinal structure of survey responses, provide interpretable latent-variable representations, and retain the familiar likelihood-based inferential framework. Moreover, political methodologists have long advocated interpreting these models through predicted probabilities and first differences rather than latent coefficients \citep{Hanmer2013a}, making them well suited for estimating the distributional causal quantities introduced in Section~\ref{subsec:estimands}.

The validity of these probability estimates, however, depends on the assumptions used to construct the conditional ordinal response distribution. Classical ordered response models impose two principal restrictions. First, the latent error distribution is assumed to belong to a fixed parametric family, typically normal or logistic. Second, the systematic component of the latent outcome is specified through a parametric single-index model. Together, these assumptions determine the shape of the conditional response distribution and therefore the estimated category probabilities.

These assumptions are often reasonable, and when approximately correct ordered response models are statistically efficient. In many survey experiments, however, they may be difficult to justify. Political attitudes are frequently polarized or clustered, producing multimodal latent distributions. Ceiling and floor effects compress responses near the boundaries of Likert scales. Persuasive treatments may affect moderate respondents differently from strong supporters, generating nonlinear treatment-response relationships. Under such conditions, misspecification of either the latent error distribution or the systematic component may distort the estimated category probabilities that form the basis of substantive interpretation.

A substantial literature has relaxed these assumptions. Generalized ordered response models allow threshold-specific effects, while semiparametric estimators relax the distributional assumption on the latent error distribution. Bayesian nonparametric approaches provide additional flexibility through hierarchical modeling. These methods represent important advances, but they typically relax only one component of the model or become computationally demanding in the presence of high-dimensional conditioning variables. The next subsection develops a unified framework that relaxes both sources of misspecification while preserving the likelihood-based structure of conventional ordered response models.

\subsection{A Flexible Likelihood-Based Framework}
\label{subsec:framework}

The preceding discussion suggests that the objective is not simply to replace the normal or logistic error distribution with a more flexible alternative. Rather, the goal is to estimate the conditional ordinal response distribution while preserving the probabilistic structure that makes ordered response models attractive. An ideal estimator should satisfy four properties. It should remain likelihood-based, accommodate high-dimensional conditioning variables, recover the full conditional response distribution, and permit varying degrees of structural modeling according to substantive knowledge.

The proposed framework retains the latent outcome representation in Equation~\eqref{eq:latent} but generalizes how its components are modeled. Recall that the latent outcome consists of two parts: the systematic component $h(\cdot)$ and the latent error distribution $F$. Existing estimators correspond to particular choices of these two components. Ordered logit and ordered probit specify both parametrically, whereas semiparametric estimators typically relax only the latent error distribution.

The proposed framework allows flexibility in either or both components while retaining a common likelihood-based estimation procedure. When substantive theory suggests that treatment and covariates operate through a common latent index, the systematic component may be specified parametrically while the latent error distribution is estimated flexibly. This specification may be viewed as a semiparametric generalization of ordered response models. Alternatively, when the functional relationship between covariates and the latent outcome is itself uncertain, the systematic component may also be estimated nonparametrically. In this way, the same estimation framework spans classical latent-index models, semiparametric ordered response models, and fully flexible conditional distribution estimators without requiring different inferential procedures.

The framework is implemented using conditional normalizing flows to model unknown probability distributions. Unlike kernel-based semiparametric estimators, conditional normalizing flows scale naturally to high-dimensional conditioning variables while retaining exact likelihood evaluation. Consequently, they preserve the familiar inferential machinery of ordered response models, including maximum-likelihood estimation, likelihood-based model comparison, and probabilistic calibration, while substantially relaxing distributional assumptions.

A particularly attractive feature of ordinal response models is that the latent error is scalar. Combined with monotone normalizing flows, this permits exact evaluation of both the latent density and its cumulative distribution function. Consequently, the threshold-crossing likelihood can be evaluated exactly without numerical integration, allowing the proposed estimator to remain computationally tractable while flexibly modeling the latent error distribution.



\subsection{Modeling the Latent Error Distribution with Normalizing Flows}
\label{subsec:flow}

% fig:nf_error {{{
\begin{figure}[ht]
\centering
\begin{tikzpicture}[>=stealth,scale=1]

% --------------------------------------------------
% Left panel: base distribution
% --------------------------------------------------
\begin{scope}[shift={(-4.6,0)}]
    \draw[thick] (-2.1,0) rectangle (2.5,3.0);
    \draw[->] (-1.75,0.25) -- (1.75,0.25) node[right] {\small $z$};
    \draw[->] (-1.75,0.25) -- (-1.75,2.35) node[above right] {\small $p_Z(z)$};

    \draw[thick]
    plot[smooth]
    coordinates{
        (-1.6,0.28)
        (-1.2,0.35)
        (-0.8,0.75)
        (-0.4,1.65)
        (0,2.25)
        (0.4,1.65)
        (0.8,0.75)
        (1.2,0.35)
        (1.6,0.28)
    };

    \node at (0,-0.35) {\small Base distribution};
    \node at (0,-0.75) {\small $z\sim N(0,1)$};
\end{scope}

% --------------------------------------------------
% Middle: conditional normalizing flow
% --------------------------------------------------
\node[draw, rounded corners, align=center, minimum width=2.8cm, minimum height=1.1cm]
    (flow) at (0,1.25)
    {\small Conditional\\[-0.1em]\small normalizing flow\\[-0.1em]\small $T_\phi(\cdot\mid D_i,X_i)$};

\draw[->,thick] (-2.25,1.35) -- (flow.west);
\draw[->,thick] (flow.east) -- (2.25,1.35);

\node[align=center] at (0,2.55)
    {\small Generative direction:\\[-0.1em]
     \small $\varepsilon_i=T_\phi(z_i\mid D_i,X_i)$};

\node[align=center] at (0,-0.55)
    {\small Normalizing direction:\\[-0.1em]
     \small $z_i=T_\phi^{-1}(\varepsilon_i\mid D_i,X_i)$};

\draw[<-,thick,dashed] (-2.25,0.65) -- (flow.west);
\draw[<-,thick,dashed] (flow.east) -- (2.25,0.65);

% --------------------------------------------------
% Right panel: flexible latent error distribution
% --------------------------------------------------
\begin{scope}[shift={(4.6,0)}]
    \draw[thick] (-2.1,0) rectangle (2.5,3.0);
    \draw[->] (-1.75,0.25) -- (1.75,0.25) node[right] {\small $\varepsilon$};
    \draw[->] (-1.75,0.25) -- (-1.75,2.35) node[above right] {\small $f_\varepsilon(\varepsilon\mid D,X)$};

    \draw[thick]
    plot[smooth]
    coordinates{
        (-1.7,0.28)
        (-1.25,0.35)
        (-0.9,0.60)
        (-0.55,0.85)
        (-0.2,0.65)
        (0.15,0.45)
        (0.45,0.75)
        (0.75,1.65)
        (1.0,2.05)
        (1.25,1.20)
        (1.55,0.35)
        (1.8,0.28)
    };

    \node at (0,-0.35) {\small Flexible latent error};
    \node at (0,-0.75) {\small distribution};
\end{scope}

% --------------------------------------------------
% Bottom: ordered response connection
% --------------------------------------------------
\node[align=center] at (0,-1.75)
{\small Combined with the latent index and thresholds, the estimated error distribution implies\\[-0.1em]
 \small conditional ordinal probabilities $P(Y_i=j\mid D_i,X_i)$.};

\end{tikzpicture}

\caption{A conditional normalizing flow transforms a simple base distribution into a flexible latent error distribution. In the proposed ordered-response framework, this replaces the fixed normal or logistic error assumption while retaining likelihood-based estimation.}
\label{fig:nf_error}
\end{figure}
% }}}

The proposed framework replaces the fixed parametric latent error distribution with a flexible conditional density estimator. Rather than assuming that the latent error follows a standard normal or logistic distribution, I estimate its conditional distribution directly from the data while retaining exact likelihood-based inference.

Normalizing flows provide a natural framework for this task. Originally introduced by \citet{Rezende2016a}, normalizing flows represent an unknown probability distribution through a sequence of invertible transformations of a simple base distribution. Because each transformation is invertible, probability densities remain analytically tractable through the change-of-variables formula, allowing estimation to proceed by maximum likelihood. Since their introduction, normalizing flows have become a standard approach for flexible likelihood-based density estimation \citep{Kobyzev2021a, Papamakarios2021a}.

Figure~\ref{fig:nf_error} illustrates the basic idea. Let

\[
z_i\sim p_Z(z)
\]

denote a simple base distribution, taken throughout this paper to be the standard normal distribution. A conditional normalizing flow transforms this base distribution into the latent error distribution through an invertible mapping,

\[
\varepsilon_i = T_\phi(z_i\mid D_i,X_i),
\]

where the transformation is parameterized by $\phi$ and conditioned on the treatment and covariates. The same transformation may be inverted,

\[
z_i = T_\phi^{-1}(\varepsilon_i\mid D_i,X_i),
\]

allowing observations to be mapped back to the base distribution. This invertibility is the defining feature of normalizing flows and enables exact likelihood evaluation through the standard change-of-variables formula,

\[
f_\varepsilon(\varepsilon_i\mid D_i,X_i) = p_Z(z_i)
\left|
\det \frac{\partial T_\phi^{-1}(\varepsilon_i\mid D_i,X_i)} {\partial\varepsilon_i}
\right|.
\]

Unlike kernel density estimators, the transformation is represented by a parametric function that scales naturally to high-dimensional conditioning variables. Throughout this paper I construct $T_\phi(\cdot)$ using monotonic rational-quadratic spline (RQS) transformations \citep{Durkan2019a}. RQS flows provide smooth monotone mappings while retaining closed-form evaluation of both the inverse transformation and the Jacobian \citep{Durkan2019a}. Because the latent error is one-dimensional, this construction also permits exact evaluation of the implied cumulative distribution function, which is required to compute the threshold probabilities in ordered response models.

The conditional flow is used only to model the latent error distribution. The systematic component of the latent outcome in Equation~\eqref{eq:latent} remains independent of this choice and may therefore be specified at different levels of structural complexity. In the structured specification, the function $h(\cdot)$ is parameterized using a conventional latent-index model, so that flexibility is introduced only through the latent error distribution. This specification may be viewed as a semiparametric generalization of ordered logit and ordered probit. In the fully flexible specification, both the systematic component and the latent error distribution are modeled nonparametrically, allowing the conditional ordinal response distribution to be learned with minimal functional-form assumptions. Both specifications share the same likelihood and estimation procedure, differing only in the amount of structure imposed on the latent outcome model.

\subsection{Likelihood-Based Estimation}
\label{subsec:likelihood}

The latent-index representation in Equation~\eqref{eq:latent}, together with the threshold-crossing measurement model introduced in Section~\ref{subsec:setup}, implies that the probability of observing response category $j$ is

\[
\begin{aligned}
P(Y_i=j\mid D_i,X_i) &= P\left( \alpha_{j-1} < Y_i^\ast \le \alpha_j \mid D_i,X_i \right) \\
&= F_\phi \left( \alpha_j - h(X_i) - D_i^\top\tau \mid D_i,X_i \right)
- F_\phi \left( \alpha_{j-1} - h(X_i) - D_i^\top\tau \mid D_i,X_i \right),
\end{aligned}
\]

where $F_\phi(\cdot)$ denotes the conditional cumulative distribution function implied by the normalizing flow introduced in the previous subsection. This expression differs from ordered logit and ordered probit only through the choice of $F_\phi$. Rather than assuming a logistic or normal distribution, the latent error distribution is estimated directly from the data. 

Also, Unlike general multivariate flow models, the one-dimensional monotone spline construction used here admits exact evaluation of the conditional cumulative distribution function $F_\phi$. Consequently, the category probabilities required by the ordered-response likelihood can be computed analytically without numerical integration.

Assuming independent observations, the log-likelihood for the observed ordinal responses is

\[
\ell(\Theta) = \sum_{i=1}^{n} \sum_{j=0}^{J} 1(Y_i=j) \log P(Y_i=j\mid D_i,X_i),
\]

where

\[
\Theta = { \tau, h, \phi, \alpha }
\]

collects all unknown parameters of the latent outcome model, the conditional normalizing flow, and the reporting thresholds.

The parameters are estimated by maximizing this likelihood using stochastic gradient optimization. Because the conditional normalizing flow admits exact evaluation of both the latent density and its cumulative distribution through the invertible transformation, the resulting objective remains a genuine likelihood rather than a surrogate prediction loss. Consequently, estimation retains many of the desirable features of conventional ordered response models, including likelihood-based model comparison and probabilistic calibration, while allowing substantially greater flexibility in the latent error distribution.

The likelihood in Equation~\eqref{eq:latent} encompasses both structured and flexible specifications discussed in the previous subsection. When the systematic component $h(\cdot)$ is parameterized linearly, the estimator may be viewed as a semiparametric generalization of ordered logit and ordered probit. When $h(\cdot)$ is itself modeled nonparametrically, the same likelihood estimates the conditional ordinal response distribution under minimal functional-form assumptions. In either case, the fitted model directly yields the conditional probabilities

\[
P(Y_i=j\mid D_i,X_i),
\]

which serve as the basis for estimating all causal quantities introduced in Section~\ref{subsec:estimands}. For example, the estimated category-specific treatment effect is obtained by standardization,

\[
\widehat{\Delta}_j = 
\frac{1}{n}
\sum_{i=1}^{n} \left[ \widehat{P}(Y_i=j\mid D_i=1,X_i) - \widehat{P}(Y_i=j\mid D_i=0,X_i)
\right],
\]

which corresponds to the plug-in estimator of the g-formula \citep{Robins1986a}. Cumulative treatment effects are computed analogously,

\[
\widehat{\Delta}_{\ge j} = 
\frac{1}{n} \sum_{i=1}^{n} \left[ \widehat{P}(Y_i\ge j\mid D_i=1,X_i) - \widehat{P}(Y_i\ge j\mid D_i=0,X_i) \right].
\]

More generally, because the entire conditional response distribution is available after estimation, any functional of that distribution may be computed without fitting additional models. Examples include average predicted probabilities, first differences, stochastic dominance measures, or global summaries of distributional change such as the one-dimensional Wasserstein (earth mover’s) distance. The proposed estimator therefore separates estimation from interpretation: the model is fitted once, while substantive quantities of interest are obtained as post-estimation summaries of the estimated response distribution.

The simulation study and empirical applications that follow focus primarily on category-specific and cumulative treatment effects because these quantities are directly interpretable and commonly reported in applied political science. The same fitted model, however, supports a substantially broader class of distributional summaries whenever alternative scientific questions arise. Throughout the simulations and empirical applications, uncertainty is quantified using the nonparametric bootstrap, which naturally propagates estimation uncertainty from the fitted conditional response distribution to all post-estimation functionals.


\clearpage

\section{Monte Carlo Simulation}
\label{sec:simulation}

This section evaluates the finite-sample performance of the proposed estimation framework. The objective is not to demonstrate that flexible estimation universally outperforms classical ordered response models. Ordered logit and ordered probit remain statistically efficient when their assumptions are approximately satisfied. Instead, the simulations investigate the conditions under which flexible estimation of the conditional ordinal response distribution improves recovery of the distributional causal effects introduced in Section~\ref{subsec:estimands}. Throughout, causal identification is held fixed so that differences in performance arise solely from the estimation procedure.

\subsection{Simulation Design}
\label{subsec:sim_design}

Each simulation begins by generating a synthetic population of one million observations from the latent-index model in Equation~\eqref{eq:latent}. The latent outcome is converted into an observed ordinal response through the threshold-crossing measurement model introduced in Section~\ref{subsec:setup}, with common reporting thresholds fixed across treatment conditions. The resulting population serves as the ground truth throughout the Monte Carlo experiment. Because the population is sufficiently large, the true category-specific and cumulative treatment effects can be computed with negligible Monte Carlo error before sampling begins. Consequently, all reported estimation error reflects finite-sample performance rather than uncertainty in the reference values.

Treatment assignment is completely randomized, ensuring that the identification assumptions developed in Section~\ref{sec:causalordinal} hold by construction. For each data-generating process, I draw 100 independent random samples of sizes $n\in{500,1000,2000}$ from the population. Each sampled dataset is analyzed independently, and performance measures are aggregated across Monte Carlo replications.

The simulation considers six data-generating processes designed to isolate different departures from the assumptions underlying classical ordered response models. Two benchmark settings generate latent errors from the normal and logistic distributions, under which ordered probit and ordered logit are correctly specified, respectively. The remaining settings introduce skewed latent errors, polarized (multimodal) latent attitudes, heteroskedastic latent errors, and nonlinear latent response functions. Together these scenarios represent empirical features frequently encountered in survey experiments while allowing the consequences of different modeling assumptions to be examined separately.

Five estimators are compared. The first is an empirical distribution estimator that estimates category probabilities directly from observed treatment-arm frequencies without imposing a statistical model. The second and third are ordered probit and ordered logit. The remaining estimators implement the proposed framework in two forms. The structured flow specification retains the linear latent-index model while estimating the latent error distribution flexibly. The model-free specification estimates both the systematic component and the latent error distribution nonparametrically. To facilitate stable optimization, both flow estimators are initialized using the ordered probit estimates before likelihood maximization.

Performance is evaluated using bias, absolute bias, and root mean squared error (RMSE) for the category-specific and cumulative treatment effects introduced in Section~\ref{subsec:estimands}. These quantities are obtained by regression standardization from each fitted model. Because all estimators target the same causal estimands under identical identification assumptions, differences in performance reflect only the statistical consequences of the modeling assumptions imposed by each estimation procedure.


\subsection{Simulation Results}
\label{subsec:sim_results}

\begin{figure}
  \begin{center}
  % \includegraphics[width=0.95\textwidth]{../r_scripts/figure_bias.pdf}
  \end{center}
  \caption{}\label{fig:bias}
\end{figure}

\begin{figure}
  \begin{center}
  % \includegraphics[width=0.95\textwidth]{../r_scripts/figure_rmse.pdf}
  \end{center}
  \caption{}\label{fig:rmse}
\end{figure}


Figures~\ref{fig:bias} and~\ref{fig:rmse} summarize the finite-sample performance of the competing estimators across the six data-generating processes. Each panel corresponds to one latent response mechanism, while the horizontal axis reports the sample size. Absolute bias and root mean squared error (RMSE) are averaged across the category-specific and cumulative treatment effects introduced in Section~\ref{subsec:estimands}. Lower values indicate more accurate recovery of the target causal quantities.

The benchmark settings provide an important validation of the proposed framework. Under normally distributed latent errors, where ordered probit is correctly specified, and under logistic latent errors, where ordered logit is correctly specified, the proposed estimators remain competitive with the classical models. In particular, the structured flow estimator exhibits little or no efficiency loss despite relaxing the parametric assumption on the latent error distribution. These benchmark scenarios therefore demonstrate that the additional modeling flexibility does not substantially compromise finite-sample performance when the conventional assumptions are approximately correct.

The remaining four scenarios examine departures from the assumptions underlying classical ordered response models. As the latent error distribution becomes skewed, multimodal, or heteroskedastic, or when the systematic component becomes nonlinear, the structured flow estimator consistently achieves the lowest estimation error among the model-based approaches. The gains are particularly pronounced under polarized latent attitudes and nonlinear response functions, where restrictive parametric assumptions lead to noticeable increases in both bias and RMSE for ordered logit and ordered probit. These results indicate that modeling the latent error distribution flexibly substantially improves recovery of the conditional ordinal response distribution when classical assumptions are violated.

The model-free flow estimator generally improves upon ordered logit and ordered probit under model misspecification but is typically outperformed by the structured flow specification. This pattern suggests that retaining the latent-index representation provides useful structural information while allowing flexibility in the latent error distribution. Conversely, the empirical distribution estimator exhibits the largest estimation error across nearly all settings because it estimates treatment-arm response probabilities directly from the observed sample without borrowing information across covariate values.

Across all simulation settings, estimation error decreases with sample size for every estimator, as expected. More importantly, the relative ranking of the estimators remains stable. The proposed framework therefore achieves the desired balance between robustness and efficiency: it remains competitive when the assumptions of classical ordered response models hold while providing substantial gains when those assumptions are violated. 


\section{Application 1: \citet{Tomz2020a}}
\label{sec:application}

The first example replicates the analysis by \citet{Tomz2020a}. The study looked into the question of how the human rights practices of foreign adversaries affect domestic public support for war through survey experiments. The main argument was that in democracies with well-established human rights, people are less willing to go to war with other countries that also respect human rights than with comparable countries that violate them. Thus they are focusing on two factors that can impact people's support for war: the political system of the foreign country and how much the foreign country respect human rights.

The paper uses total of 5 surveys, but this paper focus on the first, which contains an experiment for their main result. The survey was conducted by YouGov with a nationally representative sample of 1,430 US adults in October 2012. The treatment was given as in a form of vignettes on a country that is developing nuclear weapons. The vignettes also include trade and military relationship of the country and the US, and its conventional military strength which is set to the half of that of the US. In the vignettes, information about the political system (democracy = 1, not = 0) and human rights practices (respect = 1, not = 0) were binary treatments and they were randomized independently. 

The outcome of interest here is the support for physical strike of the US armed forces to the country. Specifically, respondents were asked to answer in 5-point scale: \textit{Favor strongly, Favor somewhat, Neither favor nor oppose, Oppose somewhat}, and \textit{Oppose strongly}. In the paper, the authors take the cardinalization approach and OLS. They first recode the outcome to take numeric values from 1 to 5, then they used this to run OLS. As discussed in Section 2, the cardinalization is arbitrary and we only can interpret the ATE as suggested by authors when the cardinalization is capturing the unknown transformation function, which is untestable.

The model they used in their main results (model 2 of their Table 1) can be expressed as following:
\begin{equation*}
    Y^{\ast} = \tau_{DEM} \cdot DEM + \tau_{HR} \cdot HR + X^{T}\beta + \varepsilon 
\end{equation*}
\begin{equation*}
    Y \in \{ 1, 2, 3, 4, 5\}
\end{equation*}

where $Y^{\ast}$ is the support for the military strike in latent scale and $Y$ is the observed outcome with cardinalization from 1 to 5, $DEM$ and $HR$ are binary variables indicating political system treatment and human rights practices treatment respectively, $X$ for the control variables and $\varepsilon$ for the error term.

Since they are interested in whether citizens living in democracy where human rights are respected (the US, to be more specific) are more reluctant to support military strike to a country who is developing nuclear weapons if the country also respecting human rights or the country has democratic political system, the causal parameters we are interested in the above equation are $\tau_{HR}$ and $\tau_{DEM}$. We only can observe the ordinal outcome, we can only identify the effects up to scale, and here I adopted error-variance normalization we discussed in the previous section. I replicate their result on Table S3 of their paper, and use ordered logit, ordered probit and KDE-based estimator as a comparison.

\begin{table}[!htb]
    \begin{center}
        \begin{tabular}{l c | c c c c }
            \\ [-1.8ex]\hline\\ [-1.8ex]
            & ATE & \multicolumn{4}{c}{NLTE} \\
            \\ [-1.8ex]\cline{2-6} \\ [-1.8ex]
            & Original -- OLS & Ordered Logit & Ordered Probit & KDE-based & NF-based \\
            \\ [-1.8ex] \hline\\ [-1.8ex]
            Respect HR & $-12.6^{\ast\ast\ast}$ & $-0.46^{\ast\ast\ast}$ & $-0.49^{\ast\ast\ast}$ & $-0.48$ & $-0.49^{\ast\ast\ast}$ \\
            & $(1.47)$ & $(0.10)$ & $(0.06)$ & $(0.27)$ & $(0.06)$ \\
            Democracy & $-9.5^{\ast\ast\ast}$ & $-0.31^{\ast\ast}$ & $-0.34^{\ast\ast\ast}$ & $-0.33$ & $-0.34^{\ast\ast\ast}$ \\
            & $(1.47)$ & $(0.10)$ & $(0.06)$ & $(0.22)$ & $(0.06)$ \\
            \\ [-1.8ex] \hline\\ [-1.8ex]
            \multicolumn{6}{l}{\scriptsize{$^{***}p<0.001$; $^{**}p<0.01$; $^{*}p<0.05$}}
        \end{tabular}
    \end{center}
    \caption{Replication of \citet{Tomz2020a} Table S3. ATE represents the original binarized OLS estimates (0-100 scale). NLTE columns report standardized latent coefficients ($\sigma=1$).}
    \label{tab:Tomz}
\end{table}

Table \ref{tab:Tomz} reports the results across all five specifications. Substantively, all models agree on the direction of the treatment effects: information that a country respects human rights or is a democracy significantly decreases domestic support for military action. However, the interpretation and statistical information utilized differ. The original OLS interpretation (a 12.6 percentage point drop for the Human Rights treatment) relies on the validity of the binarization and the cardinal 0–100 scale. By collapsing the 5-point scale into two categories, the original analysis discards potentially valuable information regarding the intensity of respondent preferences.

In contrast, NLTE models utilize the full ordinal range of the data. The standardized coefficients remain remarkably consistent across the Ordered Probit, KDE-based, and NF-based models. For instance, the effect of respecting human rights is estimated at approximately $-0.49$ standard deviations across these specifications. The fact that the flexible semiparametric estimators (KDE and NF) closely track the Ordered Probit results serves as a diagnostic, suggesting that the latent error distribution for this specific survey approximates a normal distribution. In this instance, the parametric assumption appears robust, though the semiparametric approach provides the necessary validation that the original cardinalization was not distorting the underlying direction of the effects.

\section{Application 2: \cite{Mattingly2025a}}
\label{sec:Mattingly2025a}

The second application replicates the large-scale cross-national study by \cite{Mattingly2025a}, which examines how authoritarian regimes utilize state-produced media to promote their political and economic models to foreign audiences. We focus our replication on the authors' primary hypotheses: (H1) US government messaging increases preference for the American model, (H2) Chinese Communist Party (CCP) messaging increases preference for the Chinese model, and (H3) in the presence of competing messages, public preference shifts toward the Chinese model.

To test the hypothesis, authors launched multiple survey experiments took place in 19 different countries (N = 6,000) in June 2022. Respondents were assigned via block randomization to one of four conditions: viewing US state-produced videos (USA), CCP-produced videos (China), a combination of both (Competition), or nature-themed placebo videos (Control). The primary outcomes were measured using a 6-point Likert scale, ranging from \textit{Strongly Prefer the US} (1) to \textit{Strongly Prefer China} (6). The original analysis relied on OLS regression, which cardinalizes these ordinal categories by treating the 1–6 scale as an interval measure.

After the treatment, respondents were asked to answer two questions which measures main outcome variables (whether they prefer the US (political / economic) system or (political / economic) Chinese system) in 6 categories from \textit{Strongly Prefer the US} to \textit{Strongly Prefer China}. The cardinalize the outcomes by assigning numeric values, so that \textit{Strongly Prefer the US} was recoded as 1 and \textit{Strongly Prefer China} was recoded as 6. The original results came from OLS results with these cardinalized outcomes. 

The model \cite{Mattingly2025a} used for their main results can be expressed as follows
\begin{equation*}
Y_{i}^{\ast} = \tau_{CHINA} \cdot CHINA_{i} + \tau_{USA} \cdot USA_{i} + \tau_{COMP} \cdot COMP_{i} + X_{i}^{T}\gamma + \varepsilon_{i}
\end{equation*}
\begin{equation*}
    Y \in \{ 1, 2, 3, 4, 5, 6\}
\end{equation*}

where $CHINA_{i}, USA_{i}$ and $COMP_{i}$ are indicator variables for the respective experimental conditions, and $X_{i}$ is a vector of pre-treatment covariates including age, gender, education and family income. 

While the original study estimates the Average Treatment Effect (ATE) by applying OLS directly to the categorical labels (1–6), our approach focuses on identifying the latent treatment parameters ($\tau$). As is standard in semiparametric ordinal models, the latent scale is identified only up to location and scale; hence, we apply the error-variance normalization ($\sigma = 1$). This allows for a direct comparison between the parametric specifications (Ordered Probit and Logit) and the flexible KDE and NF-based estimators. Under this framework, the coefficients represent the shift in the latent preference distribution in units of standard deviations, providing a measure of treatment intensity that is robust to the arbitrary cardinalization of the 6-point scale.

Table \ref{tab:Mattingly} reports the replication of the original OLS results, and results of NLTE estimation using different ordinal regression models including KDE-based model and flow-based model. 
For preferences over the political system (upper panel), we find a stark divergence between parametric and semiparametric models. While the direction of the effects remains consistent across all specifications, the magnitude of the estimates is significantly inflated in the cardinal (OLS) and parametric ordinal (Logit/Probit) models. Specifically, the estimated effect of the Chinese treatment in the NF-based model ($0.37$) is nearly two-thirds smaller than the OLS estimate ($1.04$) and less than half the size of the Ordered Probit estimate ($0.76$).

Notably, the \textit{Competition} effect, which the original study identified as a robust driver of preference for the Chinese model ($0.36, p < 0.001$), is substantially attenuated in the flexible models. In the NF-based specification, this effect drops to $0.11$, suggesting that when the assumption of a Normal error distribution is relaxed, the persuasive power of competitive rhetoric appears much more fragile than previously reported. This suggests that the OLS and Probit models may be over-attributing shifts in the latent distribution to the treatment, rather than accounting for the non-normal dispersion of political preferences.

In contrast, the lower panel of Table \ref{tab:Mattingly} reports preferences over economic systems, where the models show greater consensus. While OLS still yields the largest point estimates, the gap between the Probit ($0.57$) and the NF-based model ($0.58$) for the China treatment is negligible. This indicates that the latent distribution of economic preferences in this sample likely approximates normality more closely than political preferences. However, the OLS estimates for all treatments remain consistently higher than their latent counterparts, reinforcing the argument that cardinalizing Likert scales may lead to substantially different conclusions.


\begin{table}[!htb]
    \begin{center}
        \begin{tabular}{l c | c c c c }
            \\ [-1.8ex]\hline\\ [-1.8ex]
            \multicolumn{6}{l}{Outcome: Preference over Political System} \\ 
            \\ [-1.8ex] \hline\\ [-1.8ex]
            & ATE & \multicolumn{4}{c}{NLTE} \\
            \\ [-1.8ex]\cline{2-6} \\ [-1.8ex]
            & Original -- OLS & Ordered Logit & Ordered Probit & KDE-based & NF-based \\
            \\ [-1.8ex] \hline\\ [-1.8ex]
            China & $1.04^{\ast\ast\ast}$ & $0.73^{\ast\ast\ast}$ & $0.76^{\ast\ast\ast}$ & $0.36^{\ast\ast\ast}$ & $0.37^{\ast\ast\ast}$ \\
            & $(0.05)$ & $(0.04)$ & $(0.04)$ & $(0.07)$ & $(0.04)$ \\
            USA & $-0.43^{\ast\ast\ast}$ & $-0.35^{\ast\ast\ast}$ & $-0.38^{\ast\ast\ast}$ & $-0.18^{\ast\ast\ast}$ & $-0.17^{\ast\ast\ast}$ \\
            & $(0.04)$ & $(0.04)$ & $(0.04)$ & $(0.05)$ & $(0.03)$ \\
            Competition & $0.36^{\ast\ast\ast}$ & $0.21^{\ast\ast\ast}$ & $0.25^{\ast\ast\ast}$ & $0.13^{\ast}$ & $0.11^{\ast\ast}$ \\
            & $(0.05)$ & $(0.04)$ & $(0.04)$ & $(0.06)$ & $(0.04)$ \\
            \\ [-1.8ex] \hline\\ [-1.8ex]
            \multicolumn{6}{l}{Outcome: Preference over Economic System} \\ 
            \\ [-1.8ex] \hline\\ [-1.8ex]
            & ATE & \multicolumn{4}{c}{NLTE} \\
            \\ [-1.8ex]\cline{2-6} \\ [-1.8ex]
            & Original -- OLS & Ordered Logit & Ordered Probit & KDE-based & NF-based \\
            \\ [-1.8ex] \hline\\ [-1.8ex]
            China & $0.87^{\ast\ast\ast}$ & $0.56^{\ast\ast\ast}$ & $0.57^{\ast\ast\ast}$ & $0.59^{\ast\ast\ast}$ & $0.58^{\ast\ast\ast}$ \\
            & $(0.05)$ & $(0.04)$ & $(0.04)$ & $(0.07)$ & $(0.06)$ \\
            USA & $-0.57^{\ast\ast\ast}$ & $-0.39^{\ast\ast\ast}$ & $-0.43^{\ast\ast\ast}$ & $-0.42^{\ast\ast\ast}$ & $-0.43^{\ast\ast\ast}$ \\
            & $(0.05)$ & $(0.04)$ & $(0.04)$ & $(0.05)$ & $(0.05)$ \\
            Competition & $0.28^{\ast\ast\ast}$ & $0.16^{\ast\ast\ast}$ & $0.18^{\ast\ast\ast}$ & $0.19^{\ast\ast}$ & $0.18^{\ast\ast}$ \\
            & $(0.06)$ & $(0.04)$ & $(0.04)$ & $(0.06)$ & $(0.06)$ \\
            \\ [-1.8ex] \hline\\ [-1.8ex]
            \multicolumn{6}{l}{\scriptsize{$^{***}p<0.001$; $^{**}p<0.01$; $^{*}p<0.05$}}
        \end{tabular}
    \end{center}
    \caption{Replication of \citet{Mattingly2025a} Table C3. ATE represents original OLS estimates. NLTE columns report standardized latent coefficients from various ordinal specifications.}
    \label{tab:Mattingly}
\end{table}

\section{Conclusion}
\label{sec:conclusion}

This paper has examined the problem of causal inference when the outcome of interest is measured on an ordinal scale, a situation that arises frequently in political science research on attitudes and preferences. Building on a standard latent-variable framework, we showed that, even under conventional causal assumptions such as SUTVA, ignorability, and positivity, treatment effects on the underlying continuous outcome are in general identified only up to scale. The joint distribution of ordinal responses is invariant to positive affine transformations of the latent outcome and thresholds, so the absolute magnitude and location of the latent average treatment effect cannot be recovered from the data alone. Any attempt to compare coefficients (across models or across studies) must therefore adopt an explicit normalization. To address this, we proposed the \textit{Normalized Latent Treatment Effect (NLTE)} as an alternative causal estimand, ensuring that latent treatment effects are expressed in a transparent and comparable metric.

We argued that common strategies for analyzing ordinal outcomes are problematic in light of this limitation. Cardinalization, which assigns numeric scores to categories and applies tools such as OLS or difference-in-means, produces estimates whose size and sometimes even sign can depend sensitively on arbitrary labeling choices \citep{Schroder2017a, Bond2018a, Bloem2022a}. Binarization discards information about intermediate categories, leading to inefficient estimators and potentially obscuring substantively important changes in the distribution of responses. Parametric ordinal models, including ordered logit and ordered probit, avoid arbitrary numeric labels and respect the ordered nature of the data, but rely on strong assumptions about the latent error distribution. When these assumptions are violated the resulting estimators converge to pseudo-true parameters and can suffer substantial bias \citep{Manski1988a, Greene2010a, Johnston2020a, Smits2020a}.

This paper introcused two estimators that retain the interpretability of latent treatment effects while relaxing parametric assumptions on the error distribution. The first is a semiparametric KDE-based estimator that uses kernel density estimation to approximate the latent error CDF and constructs a quasi-likelihood for the ordinal model. The second employs normalizing flows to model the error distribution as an invertible transformation of a simple base distribution within a maximum likelihood framework. Both estimators treat the latent outcome as parametric, estimate the error distribution flexibly from the data.

Monte Carlo simulations demonstrated that these choices matter in practice. When the latent error distribution is correctly specified (for example, normal for ordered probit), parametric ordinal models perform well, as expected. However, under misspecification such as lognormal error, ordered logit and probit can be noticeably biased, whereas the KDE-based and NF-based estimators remain consistent and exhibit favorable finite-sample performance. The simulations also showed that cardinalization yields highly unstable estimates across different labeling schemes, and that binarization is less efficient and can miss meaningful shifts among categories. In this sense, the density-estimation-based approaches provide a robust alternative that better respects the ordinal nature of the data and the limited information it contains about the latent scale.

The empirical applications further underscore the substantive importance of these methods. Reanalyzing \citet{Tomz2020a} on human rights and military force, we found that standard binarized OLS tends to understate the magnitude of the latent treatment effect. Conversely, our replication of \citet{Mattingly2025a} regarding authoritarian propaganda reveals that the perceived "superiority" of Chinese state cues over American cues in the original study is, in part, a methodological artifact. By accounting for the floor effects and non-normal dispersion of political preferences, our proposed estimators show that the relative persuasive power of these cues is much closer than cardinal OLS suggests. These examples demonstrate how the assumption of cardinality or normality can inadvertently shape—and potentially distort—substantive conclusions in top-tier political science research.

Several directions for future work follow from these results. First, while we focused on single-item ordinal outcomes, extending density-estimation-based methods to multi-item settings and integrating them more tightly with IRT-style measurement models would be valuable. Second, developing practical diagnostic tools for assessing error distribution misspecification in ordinal models, building on surrogate residuals or related techniques, would complement the estimators proposed here. Third, applications with clustered, panel, or hierarchical data structures pose additional challenges for both identification and estimation that merit further study. More broadly, the analysis reinforces the importance of taking the ordinal nature of many political science outcomes seriously and of being explicit about the assumptions and normalizations that underlie estimates of causal effects on latent attitudes.


\bibliographystyle{apsr}
\bibliography{/Users/chanhyuk/Documents/MyLibrary}

\newpage
\section{Appendix}

\appendix

\section{KDE-Based Ordinal Regression}
\label{app:kde}

This appendix provides additional details for the semiparametric KDE-based estimator introduced in Section~\ref{sec:estimation}. We first define the quasi-likelihood, then describe the kernel estimation of the error distribution and state a consistency result.

\subsection{Model and Quasi-Likelihood}

Recall the latent outcome model
$$
Y_i^\ast = f(X_i,\beta) + D_i^\top\tau + \varepsilon_i,
$$
and the threshold-based reporting function
$$
Y_i = j
\quad\Longleftrightarrow\quad
\alpha_{j-1} < Y_i^\ast \le \alpha_j,
\qquad j = 0,1,\ldots,J,
$$
with $-\infty = \alpha_{-1} < \alpha_0 < \cdots < \alpha_J = \infty$. For notational simplicity, define the latent outcome
$$
V_i(\beta,\tau) = f(X_i,\beta) + D_i^\top\tau.
$$
Let $F$ denote the (unknown) CDF of $\varepsilon_i$. Conditional on $(X_i,D_i)$, the probability of observing category $j$ is
$$
p_j(X_i,D_i;\alpha,\beta,\tau,F)
= F\bigl( \alpha_j - V_i(\beta,\tau) \bigr)
- F\bigl( \alpha_{j-1} - V_i(\beta,\tau) \bigr).
$$
If $F$ were known, the population log-likelihood would be
$$
\ell(\alpha,\beta,\tau;F)
= \mathbb{E} \left[
\sum_{j=0}^J 1_{{Y_i=j}}
\log p_j(X_i,D_i;\alpha,\beta,\tau,F)
\right].
$$

In the sample of size $n$, the infeasible sample log-likelihood is
$$
\ell_n(\alpha,\beta,\tau;F)
= \frac{1}{n} \sum_{i=1}^n
\sum_{j=0}^J 1_{{Y_i=j}}
\log p_j(X_i,D_i;\alpha,\beta,\tau,F).
$$

When $F$ is unknown, we replace it by a nonparametric estimator $\hat{F}$ constructed via KDE, obtaining a \emph{quasi}-log-likelihood
$$
\hat{\ell}_{n}(\alpha,\beta,\tau)
= \frac{1}{n} \sum_{i=1}^n \hat{m}(X_i)
\sum_{j=0}^J 1_{{Y_i=j}}
\log \hat{p}_j(X_i,D_i;\alpha,\beta,\tau),
$$
where
$$
\hat{p}_{j}(X_i,D_i;\alpha,\beta,\tau)
= \hat{F}\bigl( \alpha_j - V_i(\beta,\tau) \bigr)
- \hat{F}\bigl( \alpha{j-1} - V_i(\beta,\tau) \bigr),
$$
and $\hat{m}(X_i)$ is a trimming function that downweights observations in extreme regions of the covariate space where KDE may be unstable. The KDE-based estimator $(\hat{\alpha},\hat{\beta},\hat{\tau})$ is defined as any maximizer of $\hat{\ell}_n(\alpha,\beta,\tau)$.

\subsection{Kernel Estimation of the Error Distribution}

To construct $\hat{F}$, we exploit the relationship between the distribution of the index $V_i$ and the error term $\varepsilon_i$. Let
$$
g_1(v \mid Y \le j) = \text{density of } V \text{ given } Y \le j,
\qquad
g_0(v \mid Y > j) = \text{density of } V \text{ given } Y > j,
$$
and let
$$
\pi_j = \mathbb{P}(Y_i \le j),
\qquad
1-\pi_j = \mathbb{P}(Y_i > j).
$$
Then by the Bayes' rule,
$$
\mathbb{P}(Y_i \le j \mid V_i = v)
= \frac{ \pi_j  g_1(v \mid Y \le j) }{
\pi_j  g_1(v \mid Y \le j)
+ (1-\pi_j)  g_0(v \mid Y > j)
}.
$$
Since $Y_i^\ast = V_i + \varepsilon_i$, and $Y_i \le j$ iff $Y_i^\ast \le \alpha_j$, we have
$$
\mathbb{P}(Y_i \le j \mid V_i = v)
= \mathbb{P}\bigl( \varepsilon_i \le \alpha_j - v \bigr)
= F\bigl( \alpha_j - v \bigr).
$$
Thus, if we can estimate $\pi_j$, $g_1(\cdot)$, and $g_0(\cdot)$ from the data, we can construct an estimator $\hat{F}$ such that
$$
\hat{F}(\alpha_j - v)
\approx \widehat{\mathbb{P}}(Y_i \le j \mid V_i = v).
$$

The probabilities $\pi_j$ and $1-\pi_j$ can be estimated by sample proportions:
$$
\hat{\pi}_{j} = \frac{1}{n} \sum_{i=1}^n 1_{{Y_i \le j}},
\qquad
1-\hat{\pi}_{j} = \frac{1}{n} \sum_{i=1}^n 1_{{Y_i > j}}.
$$
The conditional densities $g_1$ and $g_0$ are estimated using KDE. For concreteness, suppose we use a univariate kernel $K(\cdot)$ (e.g.\ Gaussian) and bandwidths that may depend on $j$. Let $n_1(j) = \sum_{i=1}^n 1_{{Y_i \le j}}$ and $n_0(j) = \sum_{i=1}^n 1_{{Y_i > j}}$. Define the subsamples
$$
{V_i : Y_i \le j}, \qquad {V_i : Y_i > j},
$$
and estimate the conditional densities as
$$
\hat{g}_{1}(v \mid Y \le j)
= \frac{1}{n_1(j) h_{1j}} \sum_{i: Y_i \le j}
K\left( \frac{v - V_i}{h_{1j}} \right),
$$
$$
\hat{g}_{0}(v \mid Y > j)
= \frac{1}{n_0(j) h_{0j}} \sum_{i: Y_i > j}
K\left( \frac{v - V_i}{h_{0j}} \right),
$$
where $h_{1j}$ and $h_{0j}$ are bandwidths chosen according to standard rules (possibly location-adaptive). Then the estimated conditional probability is
$$
\widehat{\mathbb{P}}(Y_i \le j \mid V_i = v)
= \frac{ \hat{\pi}_j  \hat{g}_1(v \mid Y \le j) }{
\hat{\pi}_j  \hat{g}_1(v \mid Y \le j)
+ (1-\hat{\pi}_j)  \hat{g}_0(v \mid Y > j)
}.
$$
We interpret this as an estimator of the CDF of $\varepsilon_i$ evaluated at $\alpha_j - v$:
$$
\hat{F}(\alpha_j - v)
\equiv
\widehat{\mathbb{P}}(Y_i \le j \mid V_i = v).
$$

In practice, the KDE is implemented with local smoothing and trimming to reduce boundary bias and instability in regions with sparse data. Following \citet{Abramson1982a} and \citet{Silverman1986a}, one can use pilot density estimates to define location-specific bandwidths and damping functions that prevent the bandwidth from shrinking too quickly in low-density regions. We adopt these techniques, as in \citet{Klein2002a}, to ensure $\sqrt{n}$ consistency.

\subsection{Consistency of the KDE-Based Estimator}

This section sketches the main consistency result for the KDE-based estimator. Fuller proofs follow from standard arguments combined with the consistency of locally smoothed KDEs \citep{Klein1993a, Abramson1982a, Silverman1986a}.

\begin{assumption}[Smoothness of the Index Distribution]
The conditional density of the index $V_i(\beta,\tau)$ given $(X_i,D_i)$ exists, is strictly positive on its support, and is continuously differentiable up to a sufficiently high order. The support of $V_i$ is compact or can be truncated to a compact set via trimming without affecting the asymptotic behavior of the estimator.
\end{assumption}

\begin{assumption}[KDE Regularity Conditions]
The kernel $K(\cdot)$ is a bounded, symmetric density with compact support. Bandwidths $h_{1j}$ and $h_{0j}$ satisfy $h_{1j} \to 0$, $h_{0j} \to 0$, and $n h_{1j} \to \infty$, $n h_{0j} \to \infty$ as $n \to \infty$, with additional conditions ensuring the consistency of locally smoothed KDEs (e.g.\ rates controlling the pilot bandwidth and damping parameters).
\end{assumption}

Under these conditions, the KDE estimators $\hat{g}_1$ and $\hat{g}_0$ converge uniformly to $g_1$ and $g_0$, and the estimated CDF $\hat{F}$ converges uniformly to the true $F$ on compact subsets of the support. Consequently, the quasi-log-likelihood $\hat{\ell}_n(\alpha,\beta,\tau)$ converges uniformly to $\ell_n(\alpha,\beta,\tau;F)$, and the maximizer of $\hat{\ell}_n$ converges to the maximizer of the infeasible log-likelihood (up to the usual positive affine transformation).

\begin{theorem}[Consistency of the KDE-Based Estimator]
Suppose the latent outcome model and threshold structure are correctly specified, and Assumptions~1--2 hold. Then, as $n \to \infty$,
$$
\bigl| (\hat{\alpha},\hat{\beta},\hat{\tau}) - (\alpha,\beta,\tau) \bigr| = o_p(1),
$$
up to an arbitrary positive affine transformation $(a + c\alpha, c\beta, c\tau)$ with $c>0$. In particular, the latent treatment effects $\tau$ are consistently estimated up to scale.
\end{theorem}

Because the ordinal model is only identified up to scale, we impose the error-variance normalization $\mathrm{Var}(\varepsilon_i)=1$ ex post by rescaling the coefficients using the estimated error variance $\hat{\sigma}_\varepsilon^2$ as described in Section~\ref{sec:estimation}. Under this normalization, the KDE-based estimator yields treatment effects that are directly comparable to those from ordered probit, ordered logit (after rescaling), and the normalizing-flow-based estimator.

\section{Normalizing Flows}
\label{app:flows}

This appendix provides additional detail on the normalizing-flow-based estimator described in Section~\ref{sec:estimation}. We briefly review the change-of-variables formula, define the ordinal likelihood with a learned error distribution, and state a consistency result.

\subsection{Change-of-Variables Formula}

Let $Z$ be a random variable with a known base density $f_Z(z)$ (for example, the standard normal density), and let $T_\theta: \mathbb{R} \to \mathbb{R}$ be an invertible, differentiable transformation parameterized by $\theta$. Define
$$
\varepsilon = T_\theta(Z).
$$
Then, by the change-of-variables formula, the density of $\varepsilon$ is
$$
f_\theta(\varepsilon)
= f_Z\bigl( T_\theta^{-1}(\varepsilon) \bigr)
\left| \frac{d}{d\varepsilon} T_\theta^{-1}(\varepsilon) \right|.
$$
The corresponding CDF $F_\theta$ is
$$
F_\theta(u)
= \int_{-\infty}^u f_\theta(t) , dt.
$$
We refer to the family ${T_\theta}_{\theta \in \Theta}$ as a \emph{normalizing flow}. By choosing a sufficiently flexible family (e.g.\ coupling layers, spline flows), we can approximate a wide class of univariate distributions with tractable densities and gradients.

\subsection{Ordinal Likelihood with a Flow-Based Error Distribution}

We embed the flow-based error distribution into the latent outcome model,
$$
Y_i^\ast = f(X_i,\beta) + D_i^\top\tau + \varepsilon_i,
\qquad
\varepsilon_i = T_\theta(Z_i),
\quad
Z_i \sim \mathcal{N}(0,1),
$$
with thresholds $\alpha$ defining the observed ordinal outcome:
$$
Y_i = j
\quad\Longleftrightarrow\quad
\alpha_{j-1} < Y_i^\ast \le \alpha_j.
$$
Let $f_\theta$ and $F_\theta$ denote the density and CDF of $\varepsilon_i$ implied by the flow $T_\theta$. Conditional on $(X_i,D_i)$, the probability of category $j$ is
$$
\mathbb{P}(Y_i = j \mid X_i,D_i;\alpha,\beta,\tau,\theta)
= F_\theta\bigl( \alpha_j - V_i(\beta,\tau) \bigr)
- F_\theta\bigl( \alpha_{j-1} - V_i(\beta,\tau) \bigr),
$$
where $V_i(\beta,\tau) = f(X_i,\beta) + D_i^\top\tau$. The sample log-likelihood is
$$
\ell_n(\alpha,\beta,\tau,\theta)
= \frac{1}{n} \sum_{i=1}^n
\sum_{j=0}^J 1_{{Y_i=j}}
\log\left[
F_\theta\bigl( \alpha_j - V_i(\beta,\tau) \bigr)
- F_\theta\bigl( \alpha_{j-1} - V_i(\beta,\tau) \bigr)
\right].
$$
We estimate $(\alpha,\beta,\tau,\theta)$ jointly by maximizing $\ell_n(\alpha,\beta,\tau,\theta)$ with respect to all parameters. Because $T_\theta$ is invertible and differentiable, both $f_\theta$ and $F_\theta$ are tractable, and gradients of the log-likelihood with respect to $(\alpha,\beta,\tau,\theta)$ can be computed efficiently using automatic differentiation.

\subsection{Consistency of the Flow-Based Estimator}

This section briefly states a consistency result for the Normalizing-Flow-Based estimator. Let $(\alpha_0,\beta_0,\tau_0,F_0)$ denote the true parameters and error distribution. Suppose that the true error distribution $F_0$ belongs to the closure of the flow family ${F_\theta : \theta \in \Theta}$, and that the latent outcome model and threshold structure are correctly specified.

\begin{assumption}[Flow Approximation and Identification]
The flow family ${F_\theta : \theta \in \Theta}$ is rich enough that there exists $\theta_0 \in \Theta$ such that $F_{\theta_0} = F_0$ (or $F_{\theta_0}$ approximates $F_0$ arbitrarily well). The parameter vector $(\alpha,\beta,\tau,\theta)$ is identified up to a positive affine transformation $(a + c\alpha, c\beta, c\tau)$ with $c>0$.
\end{assumption}

\begin{assumption}[Regularity Conditions]
The parameter space for $(\alpha,\beta,\tau,\theta)$ is compact or can be restricted to a compact subset; the log-likelihood function $\ell_n(\alpha,\beta,\tau,\theta)$ satisfies standard regularity conditions (continuity, differentiability, integrability); and the model is correctly specified in the sense that the true data-generating process is in the model class for some $(\alpha_0,\beta_0,\tau_0,\theta_0)$.
\end{assumption}

Under these conditions, standard maximum likelihood theory implies that the flow-based estimator is consistent (up to scale) and asymptotically normal.

\begin{theorem}[Consistency of the Flow-Based Estimator]
Suppose the latent outcome model and threshold structure are correctly specified, and Assumptions~3--4 hold. Then, as $n \to \infty$,
$$
\bigl| (\hat{\alpha},\hat{\beta},\hat{\tau},\hat{\theta})
- (\alpha_0,\beta_0,\tau_0,\theta_0)
\bigr| = o_p(1),
$$
up to an arbitrary positive affine transformation $(a + c\alpha_0, c\beta_0, c\tau_0)$ with $c>0$. In particular, the latent treatment effects $\tau_0$ are consistently estimated up to scale.
\end{theorem}

As with the KDE-based estimator, we impose the error-variance normalization $\mathrm{Var}(\varepsilon_i) = 1$ ex post by rescaling the coefficients using the estimated error variance. Given $\hat{\theta}$, we estimate
$$
\hat{\sigma}\varepsilon^2
= \mathrm{Var}{\hat{F}\theta}(\varepsilon_i)
\approx \frac{1}{M} \sum{m=1}^M \bigl( T_{\hat{\theta}}(Z_m) \bigr)^2,
$$
where $Z_m \sim \mathcal{N}(0,1)$ are independent draws from the base distribution. We then report
$$
\hat{\tau}^{\text{unit-var}} = \frac{\hat{\tau}}{\hat{\sigma}\varepsilon},
\qquad
\hat{\beta}^{\text{unit-var}} = \frac{\hat{\beta}}{\hat{\sigma}\varepsilon},
$$
so that all treatment effects are expressed on the same unit-error-variance latent scale as in the main text.

\end{document}



% \subsection{Transformations of Ordinal Outcomes: Cardinalization and Binarization}
%
% Researchers often attempt to circumvent the limitations of ordinal data by transforming the outcome before estimation. Two common transformations are cardinalization and binarization. We first discuss them conceptually and then use a simple example to illustrate how they can fail.
%
% By \emph{cardinalization} we mean assigning numeric labels $l_0 < \cdots < l_J$ to the ordinal categories and treating the resulting variable $\tilde{Y}_{i} = l_{{Y_i}}$ as if it were a correct representation for the latent outcome $Y_{i}^{\ast}$. Researchers then estimate treatment effects using difference-in-means, OLS, or related methods. A typical choice is $(1,2,3,4,5)$ for a five-point Likert item, but any strictly increasing set of labels would be admissible from the perspective of order.
%
% Two problems arise with cardinalization of the ordinal outcomes. First, because the labels are arbitrary, the magnitude of the estimated effect depends on the particular labeling scheme. Replacing $(1,2,3,4,5)$ with $(2,4,6,8,10)$ doubles the estimated difference-in-means without changing the underlying data. More generally, different plausible labelings can generate substantively different estimates and, in some cases, even flip the sign of the estimated effect \citep{Schroder2017a, Bond2018a, Bloem2022a}. Second, there is no guarantee that any given set of labels corresponds to the true, unknown transformation $g$ from the latent attitude to the observed categories. By putting specific numerical labels, the analyst is effectively imposing a particular form of $g$ without justification.
%
% Binarization is an coarser transformation, collapsing several ordinal categories into a single binary indicator (for example, treating \textit{strongly agree} and \textit{agree} as 1 and all other responses as 0). Although this approach has some gains from interpretation of the coefficients \citep{Breen2018a}, it discards all information about movement within the merged categories and about distinctions among intermediate categories, reducing statistical efficiency and potentially obscuring substantively meaningful shifts in attitudes. Some articles in clinical studies suggest that the use of ordinal outcomes can significantly increase the power of the study and thus reduce the required sample size by 2 to 3 times smaller compared to the binary outcomes \citep{Armstrong1989a, DAmico2020a}. Moreover, binarization cannot remedy the identification problem discussed above. It merely defines a different, coarser estimand based on coarser outcomes that may be quite distant from the underlying latent ATE.
%
% To make these issues concrete, consider a simple finite-sample example with five respondents, labeled A through E. Suppose we know each respondent’s latent potential outcomes $Y_i^\ast(0)$ and $Y_i^\ast(1)$ as in  Table~\ref{tab:latent_example}. These can be thought of as their true attitudes toward a policy under control and treatment, respectively.
%
% \begin{table}[ht]
%     \centering
%     \begin{tabular}{c|cc}
%         \hline
%         Respondent & $Y_i^\ast(0)$ & $Y_i^\ast(1)$ \\
%         \hline
%         A &  1.37 & 0.09 \\
%         B & -0.56 & 1.71 \\
%         C &  0.36 & 0.11 \\
%         D &  0.63 & 2.22 \\
%         E &  0.40 & 0.14 \\
%         \hline
%     \end{tabular}
%     \caption{Latent potential outcomes in the illustrative example.}
%     \label{tab:latent_example}
% \end{table}
%
% The true latent ATE based on Tabel~\ref{tab:latent_example} is
% $$
% \begin{aligned}
% \text{ATE}^\ast
% &= \mathbb{E}\bigl[ Y_i^\ast(1) - Y_i^\ast(0) \bigr] \\
% &= \frac{0.09 + 1.71 + 0.11 + 2.22 + 0.14}{5}\\
% &= \frac{1.37 - 0.56 + 0.36 + 0.63 + 0.40}{5} \approx 0.41. 
% \end{aligned} 
% $$ This suggests that on average, the treatment increases the latent attitude by about 0.41 units.
%
% Now suppose that we do \emph{not} observe $Y_i^\ast(d)$, but only the ordinal responses generated by applying a reporting function $g$. Consider two plausible reporting functions, $g_a$ and $g_b$, which differ in how they map latent attitudes into five ordered categories: \textit{Strongly Disagree}, \textit{Disagree}, \textit{Neither}, \textit{Agree}, and \textit{Strongly Agree}. Table~\ref{tab:g_example} shows the resulting observed potential outcomes under each $g$.
%
% \begin{table}[ht]
% \centering
% \begin{tabular}{c|cc|cc}
% \hline
% & \multicolumn{2}{c|}{$g_a$} & \multicolumn{2}{c}{$g_b$} \\
% Respondent & $Y_{i,g_a}(0)$ & $Y_{i,g_a}(1)$ & $Y_{i,g_b}(0)$ & $Y_{i,g_b}(1)$ \\
% \hline
% A & Agree          & Agree          & Agree          & Disagree \\
% B & Neither        & Agree          & Disagree       & Agree    \\
% C & Agree          & Agree          & Neither        & Disagree \\
% D & Agree          & Strongly Agree & Neither        & Agree    \\
% E & Agree          & Agree          & Neither        & Disagree \\
% \hline
% \end{tabular}
% \caption{Observed ordinal potential outcomes under two reporting functions $g_a$ and $g_b$.}
% \label{tab:g_example}
% \end{table}
%
% Both $g_a$ and $g_b$ are monotone reporting functions: higher latent values correspond to “more favorable” categories. However, they differ in how they partition the latent scale. For instance, under $g_b$ the thresholds for middle categories may be shifted relative to $g_a$, so that the same latent change produces different observed category movements.
%
% Suppose we now analyze these data using the standard cardinalization that assigns numeric scores $1,2,3,4,5$ to \textit{Strongly Disagree}, \textit{Disagree}, \textit{Neither}, \textit{Agree}, and \textit{Strongly Agree}, respectively. For each reporting function, we can compute the average cardinalized outcome under control and treatment and their difference. Under $g_a$,
% $$
% \widehat{\text{ATE}}_{\text{card},g_{a}}
% = \frac{4 + 3 + 4 + 4 + 4}{5}
% - \frac{4 + 4 + 4 + 5 + 4}{5}
% = \frac{19}{5} - \frac{21}{5}
% = -0.4,
% $$
% based on Table~\ref{tab:g_example}. Reversing the order (control minus treatment) gives $+0.4$, which happens to be close in magnitude to the true latent ATE of $0.41$. In this case, the usual $1$--$5$ coding is “lucky”: it yields a cardinalized ATE that approximates the latent ATE.
%
% Under $g_b$, applying the \emph{same} $1$--$5$ coding to the ordinal categories yields
% $$
% \widehat{\text{ATE}}_{\text{card},g_{b}}
% = \frac{4 + 2 + 3 + 3 + 3}{5}
% - \frac{2 + 4 + 2 + 4 + 2}{5}
% = \frac{15}{5} - \frac{14}{5}
% = 0.2,
% $$
% or, with the opposite ordering of potential outcomes, $-0.2$. Either way, the cardinalized ATE under $g_b$ is far from the true latent ATE of $0.41$ and can even have the wrong sign depending on whether we subtract control from treatment or vice versa. Since we do not observe $Y_i^\ast$ or know the true $g$, we cannot tell whether we are in the “lucky” case ($g_a$) or the “unlucky” case ($g_b$). This illustrates that even the standard $1$--$5$ labeling can be highly misleading, depending on the unknown reporting function.
%
% % By contrast, an ordinal regression model that correctly specifies the latent outcome (for example, $Y^\ast = \tau D + \varepsilon$) and the error distribution $F$ can, in principle, recover the underlying treatment effect up to scale using the full distribution of ordinal outcomes. Whether the thresholds are arranged as in $g_a$ or $g_b$, such a model would fit the category probabilities implied by Table~\ref{tab:g_example} and converge to a positive estimate of $\tau$ proportional to the true latent ATE. We do not develop this calculation in detail here, but the key point is that ordinal regression uses the ordering and relative frequencies of all categories rather than imposing a single arbitrary numeric scale.
%
% If we consider binarizing the outcome by defining
% $$
% Y^{\text{bin}} =
% \begin{cases}
% 1, & \text{if } Y \in {\text{Agree}, \text{Strongly Agree}}, \\
% 0, & \text{otherwise.}
% \end{cases}
% $$
% Table~\ref{tab:binary_example} reports the binarized potential outcomes and the implied average treatment effects under each reporting function.
%
% \begin{table}[ht]
% \centering
% \begin{tabular}{c|cc|cc}
% \hline
% & \multicolumn{2}{c|}{$g_a$} & \multicolumn{2}{c}{$g_b$} \\
% Respondent & $Y^{\text{bin}}{i,g_a}(0)$ & $Y^{\text{bin}}{i,g_a}(1)$ & $Y^{\text{bin}}{i,g_b}(0)$ & $Y^{\text{bin}}{i,g_b}(1)$ \\
% \hline
% A & 1 & 1 & 1 & 0 \\
% B & 0 & 1 & 0 & 1 \\
% C & 1 & 1 & 0 & 0 \\
% D & 1 & 1 & 0 & 1 \\
% E & 1 & 1 & 0 & 0 \\
% \hline
% \end{tabular}
% \caption{Binarized potential outcomes (1 if \textit{Agree} or \textit{Strongly Agree}, 0 otherwise).}
% \label{tab:binary_example}
% \end{table}
%
% Under $g_a$, the binarized means are
% $$
% \mathbb{E}[Y^{\text{bin}}{g_a}(0)] = \frac{4}{5} = 0.8,
% \qquad
% \mathbb{E}[Y^{\text{bin}}{g_a}(1)] = \frac{5}{5} = 1.0,
% $$
% so the binarized ATE is $0.2$. Under $g_b$,
% $$
% \mathbb{E}[Y^{\text{bin}}{g_b}(0)] = \frac{1}{5} = 0.2,
% \qquad
% \mathbb{E}[Y^{\text{bin}}{g_b}(1)] = \frac{2}{5} = 0.4,
% $$
% again yielding a binarized ATE of $0.2$. In both cases, binarization detects a positive effect but compresses the information in the five categories into a single bit, substantially understating the magnitude of change relative to the latent ATE of $0.41$ and failing to reflect differences between $g_a$ and $g_b$. Also, the results depend on how we collapse the ordinal outcome. If we were to choose binarize the outcomes with $1$ with \textit{Stongly Agree} and $0$ otherwise, the resulting estimates become different from what we have.
%
% % This example shows that  cardinalization with standard $1$--$5$ labels can be accidentally accurate or badly misleading depending on the unknown reporting function, an appropriately specified ordinal regression model can, in principle, recover the underlying latent effect up to scale using the full ordinal information, and binarization discards information and yields less informative and potentially inefficient estimates.
%
% \subsection{Normalization of Latent Treatment Effect}
%
% The scale and location non-identification result implies that the absolute magnitude of $Y_i^\ast$ and $\text{ATE}^\ast$ is not determined by the data. Any attempt to interpret the size of estimated coefficients or to compare them across models and samples requires a formal normalization scheme. In this section, we define our primary metric, the \textbf{Normalized Latent Treatment Effect (NLTE)}, and discuss two alternative strategies: threshold anchoring and predicted probabilities.
%
% \subsubsection{The Normalized Latent Treatment Effect (NLTE)}
% In this paper, we adopt an \emph{error-variance normalization}. Specifically, we fix the variance of the latent error term to one,
% $$
% \mathrm{Var}(\varepsilon_i) = 1.
% $$
% Under this convention, the latent outcome $Y_i^\ast$ is measured in units where the residual (unexplained) variation has standard deviation one, and each component of $\tau$ can be interpreted as the change in $Y_i^\ast$ (in these units) associated with a one-unit change in the corresponding treatment, holding other covariates fixed. This can be formally defined as below:
%
% \begin{equation}
%     \tau^{NLTE} = \frac{\mathbb{E}\left[Y^{\ast(1)} - Y^{\ast}(0)\right]}{\sqrt{Var(\varepsilon)}}
% \end{equation}
%
% This normalization aligns naturally with existing practice in ordered probit models, which impose $\varepsilon_i \sim \mathcal{N}(0,1)$ by construction. For ordered logit models, where $\varepsilon_i$ is assumed to follow a standard logistic distribution with variance $\pi^2/3$, we rescale all slope coefficients by the factor $\sqrt{3}/\pi$ to place them on the same unit-error-variance scale. That is, if $\hat{\tau}^{\text{ologit}}$ denotes the treatment coefficient from an ordered logit model, we report
% $$
% \hat{\tau}^{\text{NLTE}} = \hat{\tau}^{\text{ologit}} \times \frac{\sqrt{3}}{\pi},
% $$
% and analogously for other slope coefficients. For our density-estimation-based models, which we will introduce in the next section, we estimate the error distribution $\hat{F}$, compute its implied variance
% $$
% \hat{\sigma}_\varepsilon^2 = \mathrm{Var}_{\hat{F}}(\varepsilon_i),
% $$
% and then divide all coefficients by $\hat{\sigma}\varepsilon$:
% $$
% \hat{\tau}^{\text{NLTE}} = \frac{\hat{\tau}}{\hat{\sigma}_\varepsilon},
% $$
%
% This convention has two advantages. First, it enables \emph{direct comparison across models within the same sample}. Ordered probit, ordered logit (after rescaling), other ordinal regression models all produce treatment coefficients on the same latent scale once normalized. They measure how many units of $Y_i^\ast$ (where one unit corresponds to one standard deviation of the latent error) the treatment shifts the latent outcome. Differences in estimates can then be attributed to modeling choices (e.g.\ parametric vs.\ semiparametric error distributions) rather than arbitrary scaling. Second, it facilitates \emph{comparison across samples}, to the extent that different studies measure the same underlying construct with comparable instruments. If every model is normalized so that $\mathrm{Var}(\varepsilon_i) = 1$, then treatment coefficients from distinct datasets can be interpreted as shifts in a common latent scale on which residual variation has unit variance. Substantive differences in coefficient magnitudes across studies then reflect differences in the relationship between treatment and the latent outcome, rather than differences in ad hoc scale choices imposed by the analyst.
%
% \subsubsection{Alternative Strategies: ALTE and Probabilities}
% While the NLTE is a convenient metric for comparing latent coefficients, we also consider two alternative strategies for identification and substantive interpretation.
%
% First, we consider the \textbf{Anchored Latent Treatment Effect (ALTE)}. In case there are many treatments that can be identified, rather than fixing the variance of the unobserved error, the ALTE identifies the model by fixing the value of the effect of a treatment. Specifically, we can the value of the one coefficient as 1. Under this scheme, the error variance $\sigma_{\varepsilon}$ is a free parameter, and all other treatment coefficients are interpreted relative to the effect of the treatment serves as an anchor. 
%
% Second, we utilize \textbf{Predicted Probabilities}. To move beyond latent units entirely, we can interpret the results by calculating the change in the predicted probability of a respondent falling into a specific category $k$, or a set of categories, given a treatment $D$:
% $$
%     \Delta P = P(Y_{i} \ge k \mid D_{i} = 1, X_{i}) - P(Y_{i} \ge k \mid D_{i} = 0, X_{i})
% $$ 
% Predicted probabilities are scale-invariant and provide the most stark evidence of model divergence. Because semiparametric models (KDE and NF) can identify non-linearities and non-normalities in the error distribution, they may yield significantly different probability shifts than parametric models even if the NLTE coefficients appear similar.
%
% To clarify again, error-variance normalization does not solve the fundamental up-to-scale identification problem, and this is just one way to normalize the results. However, this provides a transparent and practically convenient metric. In the simulation and application sections that follow, we report NLTEs as our primary results, supplemented by predicted probabilities to highlight the substantive impact of relaxing parametric assumptions.
%
% \subsection{Heterogeneous Reporting Functions}
%
% The common reporting function assumption posits that all respondents use the same thresholds $(\alpha_j)$ to map their latent attitudes into ordinal responses. In practice, individuals may interpret response categories differently, leading to heterogeneity in the reporting function $g$. A simple way to incorporate such heterogeneity is to allow thresholds to depend on observed covariates $X_i$:
% $$
% Y_i = j
% \quad\Longleftrightarrow\quad
% \gamma_{j-1}^\top X_i < Y_i^\ast \le \gamma_j^\top X_i,
% \qquad j = 0,1,\ldots,J,
% $$
% where $\gamma_j$ are parameter vectors capturing how $X_i$ shifts the cutpoints. This specification is closely related to generalized ordered logit/probit and partial proportional odds models \citep{Williams2006a}.
%
% In our context, we assume that treatment indicators enter the latent outcome but not the threshold equations, so that heterogeneity in reporting is driven by background covariates rather than by the treatment itself. Under this restriction, the treatment effect $\tau$ still remains to be identified up-to-scale. Allowing for heterogeneous reporting functions can improve model fit and efficiency, and is important in applications concerned with interpersonal comparability of responses, but it does not by itself restore identification of the absolute scale of $Y_i^\ast$. We return to this specification later in the later sections.



% Political scientists has long been interested in how people form their policy preferences. In the field of international relations, determinants of domestic support for interstate conflict is one the important and widely studied topic. Recent works have found an empirical pattern between respecting human rights at home and maintaining peaceful relations abroad. That said, countries are less likely to have war if both uphold human rights compared to pairs who do not. In this vein, \cite{Tomz2020a} investigated how the human rights records of potential adversaries affect public support for war, through survey experiments in both the U.S. and the U.K.. As in usual survey experiments, individual respondents are believed to have their true opinion on using an armed forces on a unidimensional continuum, and they test if their treatment, which contains information about human rights practices of potential adversaries, can alter this opinion on average. The key dependent variable was measured with a survey question that asked respondents to answer their support for use of armed weapons against the country in 5 ordered categories: \textit{Strongly Opppose} to \textit{Strongly Favor}. This means that individual respondents should internally transformed their opinion on a continuum to one of the 5 categories.
%
% When analyzing the data, the authors took the most common approach that cardinalization of the 5 ordered categories by putting that each categories holds numeric values from 1 to 5. By taking this common practice, the authors implicitly assume that they know how individual respondents map their true continuous opinion to the 5 categories, and that mapping looks just like \textit{Strongly Oppose} means 1 and so on. However, this assumption is problematic in 2 ways. First, the decision to use the set of numeric value of 1 to 5 is completely arbitrary and can hardly justified. Second, if the assumption is violated, than the authors conclusion on the average effect of their treatment can be ill-fated. Even if one is only interested in altering the answers themselves, meaning they only intended to change the average answer from one level to the other level, this cardinalization assumption is still problematic because the change from \textit{Strongly Oppose} to \textit{Somewhat Oppose} can hardly have equal meaning as the numerical change from 1 to 2. This paper will show that we may still have some causal parameter without this strong assumption.
%
% \subsection{Identification Problem}
% % With the growing emphasis on causal inference, political science has increasingly focused on identifying the effects of policies, institutions, and political factors. Experimental and quasi-experimental methods—most prominently survey experiments, difference-in-differences, and regression discontinuity designs—have become standard empirical strategies across subfields.
% %
% % Much of this research aims to estimate causal effects on abstract and subjective outcomes that cannot be observed directly. Scholars are often interested in preferences and attitudes, such as politicians’ approval ratings \citep{Canes-Wrone2002a, Kriner2009a}, support for redistribution \citep{Alt2017a, Magni2021a}, or foreign policy orientations \citep{Scheve2001r, Mayda2005a}. These outcomes are inherently latent, yet researchers generally assume that they can be represented along a unidimensional continuum—for example, arranging individuals by the strength of their support for pension reform. To capture such constructs, survey instruments typically rely on ordered response categories, most often Likert-type items (e.g., \textit{strongly disagree, disagree, neither agree nor disagree, agree, strongly agree}).
% %
% % The challenge is that ordinal scales contain far less information than interval or cardinal measures. They indicate only the relative ordering of responses, without conveying meaningful distances between categories. Consequently, when outcomes are ordinal, researchers face the task of recovering causal effects on the unobserved latent variable using only information about order. This limitation complicates both identification and interpretation, raising important questions about how causal inference should be conducted when the outcome of interest is measured on an ordinal scale.
%
% Let us denote a binary treatment with $D$ and a latent continuous outcome variable of $Y^{\ast}$. This can be thought as people's true opinion or preference on a certain political issue or policy that may lie on a unidimensional continuous space. We can think of this as individuals internally decided support for using an armed forces in the case of \cite{Tomz2020a}. Then, inside individuals mind, this latent true opinion can be assumed to follow the partially linear true data-generating process (DGP) suggested below:
% \begin{equation}
%     Y^{\ast} = f(X, \beta) + D^{T}\tau + \varepsilon
% \end{equation}
% $X$ is a vector of factors that influence the preference or opinion other that the treatment $D$, and $\varepsilon$ is a latent error term where $\varepsilon \sim f(\cdot)$ and has a CDF of $F(\cdot)$. 
%
% When asked to express their opinion or preference in a set of ordered categories, individuals may use their internal process to transform their true latent opinion or preference, $Y^{\ast}$, to the set of ordered categories. Let's denote such internal transformation process (or \textit{reporting function}) as $g$, and the outcome of the transformation or the observed categories as $Y = \{1, 2, \ldots, J\}$. I first assume everyone shares same $g$, ignoring the additional problems arising from inter- and intra-personal differences. However, one may raise a concern that each respondent may use different internal process in transforming $Y^{\ast}$ into $Y$. This is a valid concern, I will relax this assumption later in the later Section 3.2 by making $g$ dependent on a set of covariates other than treatments analogously to \citet{Williams2006a} \footnote{This concern can further be attenuated at the design stage, by utilizing anchored vignettes suggested in \citep{King2004a, King2007a}}. 
%
% \begin{assumption}{Common Reporting Function}
%      Every respondents share a same internal process that transform their latent preferences to the observed ordinal categories.
% \end{assumption}
%
% In the usual case the categories will have labels such as \textit{Strongly Disagree / Oppose} to \textit{Strongly Favor / Agree}. We can express this relationship as:
% $$
%     Y = g(Y^{\ast}) = \begin{cases}
%         0 & \text{if } \alpha_{-1} < Y^{\ast} \le \alpha_{0} \\
%         1 & \text{if } \alpha_{0} < Y^{\ast} \le \alpha_{1} \\
%         \vdots & \vdots  \\
%         J & \text{if } \alpha_{J-1} < Y^{\ast} \le \alpha_{J} \\
%     \end{cases}
% $$, where $\alpha_{k}$ denotes the threshold points for each ordinal category. 
%
% In this setting, one natural causal estimand is the average treatment effect (ATE) in terms of $Y^{\ast}$, which equal to $\tau$ in Equation (1). Borrowing from the potential outcome framework \citep{Rubin1974a}, let us denote the potential outcome in the latent space as $Y^{\ast}(d)$, where $d \in \{0, 1\}$ denotes the treatment status. Then the ATE can be expressed as the difference between the expectations of the two potential outcomes:
% $$
%     \mathbb{E}\left[Y(1)\right] - \mathbb{E}\left[Y(0)\right]
% $$ 
%
% If we can directly observe the $Y^{\ast}$, ATE on $Y^{\ast}$ can be identified, with the following set of standard causal inference assumptions regarding the latent $Y^{\ast}$. 
%
% \begin{assumption}{SUTVA}
%     A unit's potential outcomes are not affected by the treatment given to other units. $Y^{\ast}_{i} = Y^{\ast}(D_{i})$
% \end{assumption}
%
% \begin{assumption}{Ignorability}
%     The treatment assignment is conditionally independent to the potential outcomes. $Y^{\ast}_{i} \perp\!\!\!\perp D_{i} \mid X_{i}$
% \end{assumption}
%
% \begin{assumption}{Positivity}
%     There is non-zero probability of being treated. $\mathbb{P}\left(D_i = 1\right) > 0$
% \end{assumption}
%
% The problem is that we usually \textit{do not} observe the continuous $Y^{\ast}$ directly, but only can observe the ordinal outcomes, the transformed version of $Y^{\ast}$. Thus, the naive estimator of ATE, $\mathbb{E}\left[Y^\ast(1)\right] - \mathbb{E}\left[Y^\ast(0)\right]$ cannot be obtained directly, unless strong assumptions on $g$ or $Y^{\ast}$ are imposed.
%
% One common practice to overcome this inability to have ATE is cardinalization. If we can safely assume that the latent variable $Y^{\ast}$ has numeric values that have the same length as the categories for $Y$, and that numeric values match with each category of $Y$, then we can recover $Y^{\ast}$ from $Y$ using this relationship. To clarify further, cardinalization using 1 to 5 of categories \textit{Strongly Disagree / Oppose} to \textit{Strongly Agree / Favor} can be expressed as following:
%
% \begin{assumption}{Common Cardinalization}
%     $$
%         Y = g(Y^{\ast}) = \begin{cases}
%             \text{Stronlgy Disagree} \quad & Y^\ast = 1 \\
%             \vdots \\
%             \text{Stronlgy Agree} \quad & Y^\ast = 5 \\
%         \end{cases}
%     $$ 
% \end{assumption}
%
%
% Cardinalization enables researcher to use usual causal inference tools such as means-difference, least squares, and difference-in-differences, because now the numerical values are assumed to capture the latent variable $Y^{\ast}$. One critical downside for this approach comes from the fact that the numerical labels are arbitrary, and theoretically any labels should work if they preserve the order. In some problems, researchers may find some numerical labels that can is believed to be the true values on the latent space, but in many cases it is hard to justify one labels to the other. For example, labeling \textit{strongly disagree, disagree, neither agree nor disagree, agree, strongly agree} as $\{1, 2, 3, 4, 5\}$ is no more reasonable than labeling them as $\{-1, 3, 15, 50, 100\}$. Thus, the size of the difference between expectations are hardly interpretable, since they depend on the labeling scheme. For example, in case of 5 point Likert scale, if we assign $\{2, 4, 6, 8, 10\}$, instead of $\{1, 2, 3, 4, 5\}$, the size of the difference will be doubled. This may significantly misleading, since the estimates we have only have very loose link with the treatment effects, and may over- or under-estimate them. Even worse, as discussed in \citet{Bond2018a} and \citet{Schroder2017a}, in most cases, there exists at least one labeling scheme that can flip the sign of the estimated causal effect. \citet{Bloem2022a} partly deals with this problem by providing a sensitivity test and a partial identification method based on the test. 
%
% Another problem by cardinalization is it imposes strong assumption on $g$. Therefore, if $g$ is not like in Equation (2), ATE can not be identified. To see this more clearly, suppose we have conducted same experiment as \cite{Tomz2020a} but in a much smaller scale with sample of 5. Let us further assume that we know the true potential outcome on latent space $Y^{\ast}(0)$ and $Y^{\ast}(1)$ of the respondents, and there are two candidates for $g$, $g_{a}$ and $g_{b}$ such that:
% $$
%     Y = g_{a}(Y^{\ast}) = \begin{cases}
%         \text{Strongly Disagree} &\quad -\infty < Y^\ast \le -3 \\
%         \text{Disagree} &\quad -3 < Y^\ast \le -1 \\
%         \text{Neither} &\quad -1 < Y^\ast \le 0 \\
%         \text{Agree} &\quad 0 < Y^\ast \le 2 \\
%         \text{Strongly Agree} &\quad 2 < Y^\ast < \infty \\
%     \end{cases}
% $$ 
%
% $$
%     Y = g_{b}(Y^{\ast}) = \begin{cases}
%         \text{Strongly Disagree} &\quad -\infty < Y^\ast \le -5 \\
%         \text{Disagree} &\quad -5 < Y^\ast \le 0.2 \\
%         \text{Neither} &\quad 0.2 < Y^\ast \le 1 \\
%         \text{Agree} &\quad 1 < Y^\ast \le 5 \\
%         \text{Strongly Agree} &\quad 5 < Y^\ast < \infty \\
%     \end{cases}
% $$ 
%
%     \begin{table}[ht]
%         \centering
%         \begin{tabular}{c|ll|ll|ll}
%             \hline
%             & $Y^{\ast}(0)$ & $Y^{\ast}(1)$ & $Y_{g_{a}}(0)$ & $Y_{g_{a}}(1)$ & $Y_{g_{b}}(0)$ & $Y_{g_{b}}(1)$\\
%             \hline
%             A & 1.37 & 0.09 & Agree  & Agree  & Agree  & Disagree  \\
%             B & -0.56 & 1.71  & Neither & Agree  & Disagree  & Agree \\
%             C & 0.36 & 0.11  & Agree  & Agree  & Neither & Disagree \\
%             D & 0.63 & 2.22 & Agree  & Strongly Agree & Neither & Agree  \\
%             E & 0.4 & 0.14 & Agree  & Agree  & Neither & Disagree  \\
%             \hline
%         \end{tabular}
%     \end{table}
%
% Table 1 shows the example data. The first column denote each respondents, the second and third columns show the latent potential outcomes $Y^{\ast}(0)$ and $Y^{\ast}(1)$, and the last four columns show observed potential outcomes generated by $g_{a}$ and $g_{b}$. Since we know both $Y^{\ast}(0)$ and $Y^{\ast}(1)$, we can calculate the ATE. 
% $$
%     \mathbb{E}\left[Y^\ast(1)\right] - \mathbb{E}\left[Y^\ast(0)\right] =  \frac{1.37 + (-0.56) + 0.36 + 0.63 + 0.4}{5} - \frac{0.09 + 1.71 + 0.11 + 2.22 + 0.14}{5} = 0.41
% $$ 
%
% Now suppose the data was generated based on $g_{a}$, and we put usual numeric values of 1 to 5 to each categories as in the Assumption 4. Then, the ATE based on this:
% $$
%     \frac{4 + 4 + 4 + 5 + 4}{5} - \frac{4 + 3 + 4 + 4 + 4}{5} = 0.4
% $$ 
% which is very close to the true ATE of $0.41$.
%
% However, what if the true $g$ is not $g_{a}$ but actually $g_{b}$? Then we would have observed $Y_{g_{b}}(0)$ and $Y_{g_{b}}(1)$ and the ATE calculation with the same numeric assignment of 1 to 5 as in Assumption 4:
% $$
%     \frac{4 + 2 + 3 + 3 + 3}{5} - \frac{2 + 4 + 2 + 4 + 2}{5} = -0.02
% $$ which is far off from the true ATE of $0.41$. 
%
%
%
%
% % However, the more serious problem of cardinalization approach is the interpretation of the estimates. For instance, means-difference and difference-in-differences are comparing numerical expectations of outcomes from different groups. However, if the observed outcomes are \textit{strongly disagree, disagree, neither agree nor disagree, agree, strongly agree}, even though we have transformed them into numbers, what does the expectations of these responses, mean in terms of the latent variable $Y^{\ast}$? In a potential outcome framework, the relationship between $\mathbb{E}\left[l(Y(d)) \mid X \right]$ and $\mathbb{E}\left[Y^\ast(d) \mid X \right]$ is not clear. Furthermore, even if we safely put some meanings to the expectations, 
%
% % Under this setting, we can still identify the ATE like we did in Theorem 1, if we have good knowledge on the transformation function $g$. This is because if we know $g$, then we can easily recover the latent outcome $Y^{\ast}$ form the observed outcome $Y$, by using the inverse of $g$. 
%
% % \begin{theorem}{Identification of ATE when $g$ is known}
% %     When the transformation function $g$ is known, we can identify ATE by recovering the latent outcome based on $g$ and the observed outcome.
% %     $$
% %         \mathbb{E}\left[Y^\ast(1)\right] - \mathbb{E}\left[Y^\ast(0)\right] = \mathbb{E}\left[g^{-1}(Y(1))\right] - \mathbb{E}\left[g^{-1}(Y(0))\right]
% %     $$ 
% % \end{theorem}
% %
%
% The example above demonstrate that, in case we do not observe both $Y^{\ast}$ and $g$, the identification of ATE from the common cardinalization as in Equation (2) depends almost purely on luck. 
%
%
% However, this does not mean that we should give up on causal inference in this setting as a whole. In fact, we can still identify the causal effect up to scale, based on observed cumulative probabilities and a link function. The distribution of the error term in Equation (1) is often called a link function. Link functions enable us to recover the treatment effects from the observed cumulative distribution, $P\left(Y_{i} \le j | D_{i}, X_{i}\right)$. Let us denote the distribution of $\varepsilon$ as $F$ and let $c$ be a positive constant. Using the threshold points $\alpha_{k}$, we can write ATE as:
% $$
% \begin{aligned}
%     \mathbb{E}\left[Y^{\ast}(0)\right] - \mathbb{E}\left[Y^{\ast}(1))\right] &= F^{-1}(\mathbb{P}\left(Y(1) \le j \right)) - F^{-1}(\mathbb{P}\left(Y(0) \le j \right)) \\
%     &=F^{-1} ( \mathbb{P}\left(Y^{\ast}(0) \le \alpha_j  \right) ) - F^{-1} ( \mathbb{P}\left( Y^{\ast}(1) \le  \alpha_j \right) ) \\
%     &=F_{c}^{-1}(\mathbb{P}\left(c\cdot Y^{\ast}(0) \le c \cdot \alpha_{j} \right)) - F_{c}^{-1}(\mathbb{P}\left(c \cdot Y^{\ast}(1) \le c \cdot \alpha_{j} \right)) \\
%     &=F_{c}^{-1} ( \mathbb{P}\left(c\cdot \varepsilon \le c\cdot\alpha_j - c\cdot f(X, \beta) \right) ) - F_{c}^{-1} ( \mathbb{P}\left(c\cdot\varepsilon \le c\cdot\alpha_j - c\cdot f(X, \beta) - c\cdot \tau \right) ) \\
%     &= c\cdot\tau \\
% \end{aligned}
% $$ 
% From the above, it is clear that we can identify $\tau$ up to a positive constant $c$, just because scaling up or down by $c$ would not change the observed cumulative probability. This can also be thought of as scaling up and down the transformation function $g$, which changes the numbers but will not change the observed cumulative probability $\mathbb{P}\left(Y(d) \le j\right)$.
%
% The key to the identification is now on the link function, $F$. Unfortunately, $F$ is rarely known, and we have to choose between either assuming a certain distribution or estimating it. Standard parametric ordinal regression models such as ordered logit and ordered probit models and much of the IRT models take the first approach and assume that the error term follows certain distributions. However, it is hard to justify the distributional assumption. Instead of leaning towards distributional assumptions, the second approach solve this problem by estimating the error distribution. Various distributional estimation techniques can be used here, and this paper focus on Normalizing Flows. 
%
%
% % Thus, we only can identify the size of the treatment effect relative to a positive constant $c$. Thankfully, the scaling happens at the latent variable level, so all the coefficients should be scaled by the same constant for the equation to be hold. This means that the treatment effect can be identified proportion to another coefficient. This motivates to include at least two randomized treatments in an experiments, so that we can interpret the effect of one treatment based on the other. Let us denote such treatment as an \textit{anchor}, and further denote $\tau_{a}$ as a coefficient of the anchor, and we can define the anchored latent treatment effect as:
% %
% % \begin{definition}{Anchored Latent Treatment Effect}
% %     $$
% %         \tau_{ALTE} = \frac{c\cdot \tau}{c \cdot \tau_{a}} =  \frac{\tau}{\tau_{a}}
% %     $$
% % \end{definition}
%
% % To our running example, \citet{Tomz2020a} randomized both the political system of the adversaries (democracy = 1 or not = 0) and their human rights practice (respect = 1, not = 0). Both treatment effects can only be identified up to a common constant, thus choosing one of them as an anchor will make the other treatment more interpretable. For example, if we choose to use political system as an anchor, then the ALTE of the human rights practice treatment can be interpreted as how much bigger or smaller relative to the treatment effect of the political system. 
%
% % {\color{blue} \textit{Potential alternatives}
% %
% % a) We may measure the treatment effect in terms of the standard deviation of the given data. That is, we can normalize the coefficients based on the standard deviation of all covariates. Let $X' = X + D$ than, we can standardize the coefficients based on the the standard deviation of the data, $sd(X')$. 
% %
% % b) we may use $sd(\varepsilon)$ as a normalizing constant. But can be tricky with KDE.
% %
% % c) Focus only on distributional changes.
% % $$
% %     DTE_{j} = P(Y = j \mid D = 1)  - P(Y = j \mid D = 0)
% % $$ 
% % Downside is that it is hard to average the effect; without proper weights for each category $j$, it may bias.
% %
% % d) if we use normalizing flows instead of KDE, we may scale coefficients to the standard normal distribution... not 100\% sure.
% %
% % }
%
%
%
% \section{Estimation}
%
% In this section, we visit two common distributional assumption based approaches for up to scale identification and introduce distributional assumption based approaches. 
%
% \subsection{Two Dominant Approaches with Distributional Assumptions}
%
% IRT based approaches and standard ordered logit and probit regression are the most two common methods that can estimate the causal effect with ordinal outcomes. IRT based on two-step approach. First, it estimates the latent variable using the series of items that were designed to measure the same concept in different angles. In the context of the model setting above, using multiple $Y$, IRT estimates the $Y^{\ast}$. After we get the estimate for the latent variable $Y^{\ast}$, then we can employ the usual causal inference tools to identify $\tau$ up to scale since it is now have numerical meanings. The second variant of IRT based approach, the hierarchical IRT is recent discussed in \citet{Stoetzer2025a}. Instead of estimating the latent variable in the first step, they effectively merge the two steps in to EM algorithm, and produces more consistent estimates of the causal effects. Since IRT based methods do not put any numerical meaning to the ordinal outcomes themselves, the estimates from the methods can be interpreted as the target causal effect up to scale. One downside of IRT approach is that it is advised to have at least 3 different items that are assumed to measure the same concept. However, most of the political science research rely on 1 or 2 items in measuring their outcomes. Therefore, if researcher cares the treatment effect on specific 1 or 2 items, IRT might not be a good choice. Furthermore, many IRT approaches shares the same shortcoming of assuming the distribution of $F$ as standard ordered logit and probit regressions, which I will elaborate more below.
%
% Another approach is using ordinal regression, and by far the most common methods are ordered logit and probit regressions. Similar to the IRT based methods, both ordered logit and probit utilize the ordered nature of the outcomes measured in ordinal scales, and designed to identify the actual target causal effect in unobserved, latent space up to scale. However, one big shortcoming of most of the IRT models and standard ordinal regression models is that they require strong distributional assumptions on the error term in the latent space, $\varepsilon$.
%
% \begin{assumption}{Distributional Assumption for IRT and Ordinal Regression Models}
%     The error term, $\varepsilon$ follows a certain distribution. For example, standard logistic distribution (ordered logit model) or standard normal distribution (ordered probit model). 
% \end{assumption}
%
% While these assumptions facilitate fast and efficient estimation through maximum likelihood, when the distributional assumptions fail, the estimates are statistically inconsistent and biased. 
%
% Let $i$ be the index for $i$th data. To construct likelihood function, the probability of observing $Y = j$ given $X_{i}$ is:
% $$
% \begin{aligned}
%     \mathbb{P}\left(Y = j \mid X_{i}\right) &= \mathbb{P}\left(Y^\ast \le \alpha_j\right) - \mathbb{P}\left(Y^\ast > \alpha_{j-1}\right) \\
%     &=\mathbb{P}\left(\varepsilon \le \alpha_j - (f(X_{i}, \beta) + D_{i}^{T}\tau) \right) - \mathbb{P}\left(\varepsilon > \alpha_{j-1} - (f(X_{i}, \beta) + D_{i}^{T}\tau) \right)\\
%     &= F(\alpha_{j} - (f(X_{i}, \beta) + D_{i}^{T}\tau) ) - F(\alpha_{j-1} - (f(X_{i}, \beta) + D_{i}^{T}\tau))
% \end{aligned}
% $$
% , where $F(\cdot)$ is a unknown but assumed distributional function (CDF). 
%
% Based on this, the parameters can be easily estimated with maximum likelihood.
% \begin{equation}
%      (\beta, \tau) = \arg\max_{\beta, \tau} \mathbb{E}\left[\frac{1}{n} \sum_{i=1}^{n} \log \left( F(\alpha_{j} - (f(X_{i}, \beta) + D_{i}^{T}\tau)) - F(\alpha_{j-1} - (f(X_{i}, \beta) + D_{i}^{T}\tau)) \right) \right]
% \end{equation}
%
% To solve this MLE, standard ordered logit and probit models put assumption on $F(\cdot)$. For example, Ordered Probit model assume that the error term ($\varepsilon$) follows the standard normal distribution. Let $\Phi(\cdot)$ denote the standard normal distribution function. Then the MLE becomes:
%
% $$
%     (\beta, \tau) = \arg\max_{\beta, \tau} \mathbb{E}\left[\frac{1}{n} \sum_{i=1}^{n} \log \left( \Phi(\alpha_{j} - (f(X_{i}, \beta) + D_{i}^{T}\tau)) - \Phi(\alpha_{j-1} - (f(X_{i}, \beta) + D_{i}^{T}\tau)) \right) \right]
% $$ 
%
% However, it is not guaranteed that the distributional assumption that $F(\cdot) = \Phi(\cdot)$, and if the assumption is violated, the estimators from ordered probit model are to be biased and inconsistent, because it will converge to some value $(\tilde{\beta}, \tilde{\alpha}) \neq (\beta, \alpha)$. Increase in sample size does not make the two distributions closer, and therefore the estimators from two MLE will not converge to the true values making it inconsistent \citep{Manski1988a, Greene2010a, Bond2018a}. Considering ordered logit model, logistic distribution is assumed and exactly same logic will lead to the biased and inconsistent estimation of parameters. The size and direction of the bias depend on the difference between $F(\cdot)$ and the CDF of the assumed distribution, cannot expect precisely, but in recent empirical works, right (positively) skewed error distributions would generally attenuate the size of the $\beta$ estimations \citep{Johnston2020a, Smits2020a}. This issue even harder to detect because ordered logit and ordered probit models often produce similar results. 
%
% % Theoretically, the estimates from a misspecified model may have opposite sign to the true treatment effect, leading researchers toward substantially different inferences \citep{Manski1988a, Greene2010a}. 
%
% As demonstrated, both IRT and standard parametric ordinal regressions rely on strong assumptions. The distributional assumption is especially concerning, since this can lead to inconsistent and biased estimation of the coefficients. To overcome this shortcoming, this paper introduces alternative modeling based on Kernel density estimation (KDE) and Normalizing flows. Both models relax the distributional assumption and guarantees consistent estimation of coefficients by estimating the error distribution from the given data instead of assuming certain distributions.
